% =====================================================================
%  Anatomía de un LLM — Fase 0 — Módulo 0.6
%  Cálculo integral y series
%  Compilar: latexmk -r ../latexmkrc P06-calculo-integral.tex
% =====================================================================
\documentclass[11pt,a4paper,twoside,openany]{book}
\usepackage{../estilo-apuntes}

\newcommand{\ApuntesCurso}{Anatomía de un LLM}
\newcommand{\ApuntesModulo}{Módulo 0.6}

% Notación propia del módulo
\newcommand{\N}{\mathbb{N}}
\newcommand{\Z}{\mathbb{Z}}
\newcommand{\Q}{\mathbb{Q}}
\newcommand{\dif}{\mathop{}\!d}
\newcommand{\Lsum}{\operatorname{L}}
\newcommand{\Usum}{\operatorname{U}}
\newcommand{\Rsum}{\operatorname{S}}
\newcommand{\eval}[3]{\bigl[#1\bigr]_{#2}^{#3}}

% Anclas de hiperenlace únicas para los enunciados (el contador se reinicia en cada tema)
\makeatletter
\@for\@env:={definicion,teorema,proposicion,lema,corolario,ejemplo,observacion}\do{%
  \expandafter\providecommand\csname theH\@env\endcsname{}%
  \expandafter\renewcommand\csname theH\@env\endcsname{\thechapter.\arabic{definicion}}}
\makeatother

\hypersetup{pdftitle={Módulo 0.6. Cálculo integral y series (resumen)},pdfauthor={Anatomía de un LLM}}

\makeatletter
\newlabel{ej:primitiva-condicion}{{1.3}{1}{}{ejemplo.1.3}{}}
\newlabel{ej:primitiva-condicion@cref}{{[ejemplo][3][]1.3}{[1][1][]1}}
\newlabel{fig:darboux}{{1.1}{3}{Sumas de Darboux de $f(x)=x^{2}$ en $[0,1]$ con la partición uniforme de cuatro subintervalos. La suma inferior es el área gris oscuro, $\Lsum =\frac {14}{64}$; la superior añade el gris claro, $\Usum =\frac {30}{64}$. El valor exacto, $\frac 13$, queda entre ambas}{figure.1.1}{}}
\newlabel{fig:darboux@cref}{{[figure][1][1]1.1}{[1][2][]3}}
\newlabel{ej:x2}{{1.10}{3}{Integral de $x^{2}$}{ejemplo.1.10}{}}
\newlabel{ej:x2@cref}{{[ejemplo][10][]1.10}{[1][3][]3}}
\newlabel{ej:dirichlet}{{1.11}{4}{Una función no integrable}{ejemplo.1.11}{}}
\newlabel{ej:dirichlet@cref}{{[ejemplo][11][]1.11}{[1][4][]4}}
\newlabel{ej:estimacion}{{1.17}{5}{Estimación de una integral}{ejemplo.1.17}{}}
\newlabel{ej:estimacion@cref}{{[ejemplo][17][]1.17}{[1][5][]5}}
\newlabel{obs:1-errores}{{1.18}{5}{Errores frecuentes}{observacion.1.18}{}}
\newlabel{obs:1-errores@cref}{{[observacion][18][]1.18}{[1][5][]5}}
\newlabel{eje:1-inmediatas}{{1.1}{6}{}{ejercicio.1.1}{}}
\newlabel{eje:1-inmediatas@cref}{{[ejercicio][1][1]1.1}{[1][6][]6}}
\newlabel{eje:1-vf}{{1.2}{6}{}{ejercicio.1.2}{}}
\newlabel{eje:1-vf@cref}{{[ejercicio][2][1]1.2}{[1][6][]6}}
\newlabel{eje:1-sumas-x3}{{1.3}{6}{}{ejercicio.1.3}{}}
\newlabel{eje:1-sumas-x3@cref}{{[ejercicio][3][1]1.3}{[1][6][]6}}
\newlabel{eje:1-sumas-strang}{{1.4}{6}{}{ejercicio.1.4}{}}
\newlabel{eje:1-sumas-strang@cref}{{[ejercicio][4][1]1.4}{[1][6][]6}}
\newlabel{eje:1-limite-x3}{{1.5}{6}{}{ejercicio.1.5}{}}
\newlabel{eje:1-limite-x3@cref}{{[ejercicio][5][1]1.5}{[1][6][]6}}
\newlabel{eje:1-limite-seno}{{1.6}{6}{}{ejercicio.1.6}{}}
\newlabel{eje:1-limite-seno@cref}{{[ejercicio][6][1]1.6}{[1][6][]6}}
\newlabel{eje:1-limites-sumas}{{1.7}{6}{}{ejercicio.1.7}{}}
\newlabel{eje:1-limites-sumas@cref}{{[ejercicio][7][1]1.7}{[1][6][]6}}
\newlabel{eje:1-condicion}{{1.8}{6}{}{ejercicio.1.8}{}}
\newlabel{eje:1-condicion@cref}{{[ejercicio][8][1]1.8}{[1][6][]6}}
\newlabel{eje:1-estimar}{{1.9}{7}{}{ejercicio.1.9}{}}
\newlabel{eje:1-estimar@cref}{{[ejercicio][9][1]1.9}{[1][6][]7}}
\newlabel{eje:1-exp-sumas}{{1.10}{7}{}{ejercicio.1.10}{}}
\newlabel{eje:1-exp-sumas@cref}{{[ejercicio][10][1]1.10}{[1][7][]7}}
\newlabel{eje:1-diferencia}{{1.11}{7}{}{ejercicio.1.11}{}}
\newlabel{eje:1-diferencia@cref}{{[ejercicio][11][1]1.11}{[1][7][]7}}
\newlabel{eje:1-positiva}{{1.12}{7}{}{ejercicio.1.12}{}}
\newlabel{eje:1-positiva@cref}{{[ejercicio][12][1]1.12}{[1][7][]7}}
\newlabel{eje:1-un-punto}{{1.13}{7}{}{ejercicio.1.13}{}}
\newlabel{eje:1-un-punto@cref}{{[ejercicio][13][1]1.13}{[1][7][]7}}
\newlabel{eje:1-nula}{{1.14}{7}{}{ejercicio.1.14}{}}
\newlabel{eje:1-nula@cref}{{[ejercicio][14][1]1.14}{[1][7][]7}}
\newlabel{eje:1-racionales-irr}{{1.15}{7}{}{ejercicio.1.15}{}}
\newlabel{eje:1-racionales-irr@cref}{{[ejercicio][15][1]1.15}{[1][7][]7}}
\newlabel{ej:barrow}{{2.7}{9}{}{ejemplo.2.7}{}}
\newlabel{ej:barrow@cref}{{[ejemplo][7][]2.7}{[1][9][]9}}
\newlabel{ej:area-parabola}{{2.10}{10}{}{ejemplo.2.10}{}}
\newlabel{ej:area-parabola@cref}{{[ejemplo][10][]2.10}{[1][10][]10}}
\newlabel{ej:area-seno}{{2.11}{10}{}{ejemplo.2.11}{}}
\newlabel{ej:area-seno@cref}{{[ejemplo][11][]2.11}{[1][10][]10}}
\newlabel{ej:media-seno}{{2.12}{10}{}{ejemplo.2.12}{}}
\newlabel{ej:media-seno@cref}{{[ejemplo][12][]2.12}{[1][10][]10}}
\newlabel{obs:2-errores}{{2.15}{11}{Errores frecuentes}{observacion.2.15}{}}
\newlabel{obs:2-errores@cref}{{[observacion][15][]2.15}{[1][11][]11}}
\newlabel{eje:2-conceptos}{{2.1}{11}{}{ejercicio.2.1}{}}
\newlabel{eje:2-conceptos@cref}{{[ejercicio][1][2]2.1}{[1][11][]11}}
\newlabel{eje:2-error-barrow}{{2.2}{11}{}{ejercicio.2.2}{}}
\newlabel{eje:2-error-barrow@cref}{{[ejercicio][2][2]2.2}{[1][11][]11}}
\newlabel{eje:2-barrow}{{2.3}{11}{}{ejercicio.2.3}{}}
\newlabel{eje:2-barrow@cref}{{[ejercicio][3][2]2.3}{[1][11][]11}}
\newlabel{eje:2-arco-seno}{{2.4}{11}{}{ejercicio.2.4}{}}
\newlabel{eje:2-arco-seno@cref}{{[ejercicio][4][2]2.4}{[1][11][]11}}
\newlabel{eje:2-parabola-eje}{{2.5}{12}{}{ejercicio.2.5}{}}
\newlabel{eje:2-parabola-eje@cref}{{[ejercicio][5][2]2.5}{[1][12][]12}}
\newlabel{eje:2-derivar}{{2.6}{12}{}{ejercicio.2.6}{}}
\newlabel{eje:2-derivar@cref}{{[ejercicio][6][2]2.6}{[1][12][]12}}
\newlabel{eje:2-hallar-f}{{2.7}{12}{}{ejercicio.2.7}{}}
\newlabel{eje:2-hallar-f@cref}{{[ejercicio][7][2]2.7}{[1][12][]12}}
\newlabel{eje:2-valor-medio}{{2.8}{12}{}{ejercicio.2.8}{}}
\newlabel{eje:2-valor-medio@cref}{{[ejercicio][8][2]2.8}{[1][12][]12}}
\newlabel{eje:2-estudio-F}{{2.9}{12}{}{ejercicio.2.9}{}}
\newlabel{eje:2-estudio-F@cref}{{[ejercicio][9][2]2.9}{[1][12][]12}}
\newlabel{eje:2-limite}{{2.10}{12}{}{ejercicio.2.10}{}}
\newlabel{eje:2-limite@cref}{{[ejercicio][10][2]2.10}{[1][12][]12}}
\newlabel{eje:2-par-impar}{{2.11}{12}{}{ejercicio.2.11}{}}
\newlabel{eje:2-par-impar@cref}{{[ejercicio][11][2]2.11}{[1][12][]12}}
\newlabel{eje:2-logaritmo}{{2.12}{12}{}{ejercicio.2.12}{}}
\newlabel{eje:2-logaritmo@cref}{{[ejercicio][12][2]2.12}{[1][12][]12}}
\newlabel{eje:2-periodica}{{2.13}{12}{}{ejercicio.2.13}{}}
\newlabel{eje:2-periodica@cref}{{[ejercicio][13][2]2.13}{[1][12][]12}}
\newlabel{eje:2-simetrica-f}{{2.14}{12}{}{ejercicio.2.14}{}}
\newlabel{eje:2-simetrica-f@cref}{{[ejercicio][14][2]2.14}{[1][12][]12}}
\newlabel{eje:2-tvm-ponderado}{{2.15}{13}{}{ejercicio.2.15}{}}
\newlabel{eje:2-tvm-ponderado@cref}{{[ejercicio][15][2]2.15}{[1][12][]13}}
\newlabel{ej:sustitucion}{{3.2}{14}{}{ejemplo.3.2}{}}
\newlabel{ej:sustitucion@cref}{{[ejemplo][2][]3.2}{[1][14][]14}}
\newlabel{ej:circulo}{{3.3}{14}{Área del círculo}{ejemplo.3.3}{}}
\newlabel{ej:circulo@cref}{{[ejemplo][3][]3.3}{[1][14][]14}}
\newlabel{ej:partes}{{3.5}{15}{}{ejemplo.3.5}{}}
\newlabel{ej:partes@cref}{{[ejemplo][5][]3.5}{[1][15][]15}}
\newlabel{ej:fracciones}{{3.8}{16}{}{ejemplo.3.8}{}}
\newlabel{ej:fracciones@cref}{{[ejemplo][8][]3.8}{[1][16][]16}}
\newlabel{ej:trig}{{3.9}{16}{}{ejemplo.3.9}{}}
\newlabel{ej:trig@cref}{{[ejemplo][9][]3.9}{[1][16][]16}}
\newlabel{obs:3-errores}{{3.10}{17}{Errores frecuentes}{observacion.3.10}{}}
\newlabel{obs:3-errores@cref}{{[observacion][10][]3.10}{[1][16][]17}}
\newlabel{eje:3-elegir}{{3.1}{17}{}{ejercicio.3.1}{}}
\newlabel{eje:3-elegir@cref}{{[ejercicio][1][3]3.1}{[1][17][]17}}
\newlabel{eje:3-sustitucion}{{3.2}{17}{}{ejercicio.3.2}{}}
\newlabel{eje:3-sustitucion@cref}{{[ejercicio][2][3]3.2}{[1][17][]17}}
\newlabel{eje:3-limites}{{3.3}{17}{}{ejercicio.3.3}{}}
\newlabel{eje:3-limites@cref}{{[ejercicio][3][3]3.3}{[1][17][]17}}
\newlabel{eje:3-partes}{{3.4}{17}{}{ejercicio.3.4}{}}
\newlabel{eje:3-partes@cref}{{[ejercicio][4][3]3.4}{[1][17][]17}}
\newlabel{eje:3-ciclicas}{{3.5}{17}{}{ejercicio.3.5}{}}
\newlabel{eje:3-ciclicas@cref}{{[ejercicio][5][3]3.5}{[1][17][]17}}
\newlabel{eje:3-fracciones}{{3.6}{18}{}{ejercicio.3.6}{}}
\newlabel{eje:3-fracciones@cref}{{[ejercicio][6][3]3.6}{[1][17][]18}}
\newlabel{eje:3-fracciones-cuad}{{3.7}{18}{}{ejercicio.3.7}{}}
\newlabel{eje:3-fracciones-cuad@cref}{{[ejercicio][7][3]3.7}{[1][18][]18}}
\newlabel{eje:3-racional-impropia}{{3.8}{18}{}{ejercicio.3.8}{}}
\newlabel{eje:3-racional-impropia@cref}{{[ejercicio][8][3]3.8}{[1][18][]18}}
\newlabel{eje:3-trig}{{3.9}{18}{}{ejercicio.3.9}{}}
\newlabel{eje:3-trig@cref}{{[ejercicio][9][3]3.9}{[1][18][]18}}
\newlabel{eje:3-trig-sust}{{3.10}{18}{}{ejercicio.3.10}{}}
\newlabel{eje:3-trig-sust@cref}{{[ejercicio][10][3]3.10}{[1][18][]18}}
\newlabel{eje:3-curso}{{3.11}{18}{}{ejercicio.3.11}{}}
\newlabel{eje:3-curso@cref}{{[ejercicio][11][3]3.11}{[1][18][]18}}
\newlabel{eje:3-reduccion}{{3.12}{18}{}{ejercicio.3.12}{}}
\newlabel{eje:3-reduccion@cref}{{[ejercicio][12][3]3.12}{[1][18][]18}}
\newlabel{eje:3-wallis}{{3.13}{18}{}{ejercicio.3.13}{}}
\newlabel{eje:3-wallis@cref}{{[ejercicio][13][3]3.13}{[1][18][]18}}
\newlabel{eje:3-beta}{{3.14}{18}{}{ejercicio.3.14}{}}
\newlabel{eje:3-beta@cref}{{[ejercicio][14][3]3.14}{[1][18][]18}}
\newlabel{eje:3-sustitucion-falsa}{{3.15}{18}{}{ejercicio.3.15}{}}
\newlabel{eje:3-sustitucion-falsa@cref}{{[ejercicio][15][3]3.15}{[1][18][]18}}
\newlabel{ej:impropias1}{{4.3}{19}{}{ejemplo.4.3}{}}
\newlabel{ej:impropias1@cref}{{[ejemplo][3][]4.3}{[1][19][]19}}
\newlabel{ej:log-impropia}{{4.6}{20}{}{ejemplo.4.6}{}}
\newlabel{ej:log-impropia@cref}{{[ejemplo][6][]4.6}{[1][20][]20}}
\newlabel{ej:criterios}{{4.10}{20}{}{ejemplo.4.10}{}}
\newlabel{ej:criterios@cref}{{[ejemplo][10][]4.10}{[1][20][]20}}
\newlabel{obs:4-errores}{{4.16}{22}{Errores frecuentes}{observacion.4.16}{}}
\newlabel{obs:4-errores@cref}{{[observacion][16][]4.16}{[1][22][]22}}
\newlabel{eje:4-clasificar}{{4.1}{22}{}{ejercicio.4.1}{}}
\newlabel{eje:4-clasificar@cref}{{[ejercicio][1][4]4.1}{[1][22][]22}}
\newlabel{eje:4-calcular}{{4.2}{22}{}{ejercicio.4.2}{}}
\newlabel{eje:4-calcular@cref}{{[ejercicio][2][4]4.2}{[1][22][]22}}
\newlabel{eje:4-logaritmos}{{4.3}{23}{}{ejercicio.4.3}{}}
\newlabel{eje:4-logaritmos@cref}{{[ejercicio][3][4]4.3}{[1][22][]23}}
\newlabel{eje:4-comparar}{{4.4}{23}{}{ejercicio.4.4}{}}
\newlabel{eje:4-comparar@cref}{{[ejercicio][4][4]4.4}{[1][23][]23}}
\newlabel{eje:4-seno}{{4.5}{23}{}{ejercicio.4.5}{}}
\newlabel{eje:4-seno@cref}{{[ejercicio][5][4]4.5}{[1][23][]23}}
\newlabel{eje:4-gamma-valores}{{4.6}{23}{}{ejercicio.4.6}{}}
\newlabel{eje:4-gamma-valores@cref}{{[ejercicio][6][4]4.6}{[1][23][]23}}
\newlabel{eje:4-gamma-medio}{{4.7}{23}{}{ejercicio.4.7}{}}
\newlabel{eje:4-gamma-medio@cref}{{[ejercicio][7][4]4.7}{[1][23][]23}}
\newlabel{eje:4-erf}{{4.8}{23}{}{ejercicio.4.8}{}}
\newlabel{eje:4-erf@cref}{{[ejercicio][8][4]4.8}{[1][23][]23}}
\newlabel{eje:4-normal}{{4.9}{23}{}{ejercicio.4.9}{}}
\newlabel{eje:4-normal@cref}{{[ejercicio][9][4]4.9}{[1][23][]23}}
\newlabel{eje:4-gamma-conv}{{4.10}{23}{}{ejercicio.4.10}{}}
\newlabel{eje:4-gamma-conv@cref}{{[ejercicio][10][4]4.10}{[1][23][]23}}
\newlabel{eje:4-limite-comp}{{4.11}{23}{}{ejercicio.4.11}{}}
\newlabel{eje:4-limite-comp@cref}{{[ejercicio][11][4]4.11}{[1][23][]23}}
\newlabel{eje:4-dirichlet}{{4.12}{23}{}{ejercicio.4.12}{}}
\newlabel{eje:4-dirichlet@cref}{{[ejercicio][12][4]4.12}{[1][23][]23}}
\newlabel{eje:4-potencias}{{4.13}{23}{}{ejercicio.4.13}{}}
\newlabel{eje:4-potencias@cref}{{[ejercicio][13][4]4.13}{[1][23][]23}}
\newlabel{eje:4-decreciente}{{4.14}{23}{}{ejercicio.4.14}{}}
\newlabel{eje:4-decreciente@cref}{{[ejercicio][14][4]4.14}{[1][23][]23}}
\newlabel{eje:4-wallis-gauss}{{4.15}{24}{}{ejercicio.4.15}{}}
\newlabel{eje:4-wallis-gauss@cref}{{[ejercicio][15][4]4.15}{[1][23][]24}}
\newlabel{ej:potencias}{{5.3}{25}{Potencias}{ejemplo.5.3}{}}
\newlabel{ej:potencias@cref}{{[ejemplo][3][]5.3}{[1][25][]25}}
\newlabel{ej:telescopica}{{5.8}{26}{Series telescópicas}{ejemplo.5.8}{}}
\newlabel{ej:telescopica@cref}{{[ejemplo][8][]5.8}{[1][26][]26}}
\newlabel{fig:criterio-integral}{{5.1}{27}{Criterio de la integral con $f(x)=1/x$. El rectángulo sobre $[k,k+1]$ tiene altura $a_k=f(k)$ y contiene la región bajo la curva en ese intervalo; por eso $\sum _{k=1}^{n}a_k\ge \int _1^{n+1}f$. Desplazando los rectángulos una unidad a la izquierda quedan por debajo de la curva, y se obtiene la otra desigualdad}{figure.5.1}{}}
\newlabel{fig:criterio-integral@cref}{{[figure][1][5]5.1}{[1][27][]27}}
\newlabel{ej:cociente}{{5.15}{28}{}{ejemplo.5.15}{}}
\newlabel{ej:cociente@cref}{{[ejemplo][15][]5.15}{[1][28][]28}}
\newlabel{ej:armonica-alternada}{{5.19}{28}{Serie armónica alternada}{ejemplo.5.19}{}}
\newlabel{ej:armonica-alternada@cref}{{[ejemplo][19][]5.19}{[1][28][]28}}
\newlabel{obs:5-errores}{{5.21}{28}{Errores frecuentes}{observacion.5.21}{}}
\newlabel{obs:5-errores@cref}{{[observacion][21][]5.21}{[1][28][]28}}
\newlabel{eje:5-vf}{{5.1}{29}{}{ejercicio.5.1}{}}
\newlabel{eje:5-vf@cref}{{[ejercicio][1][5]5.1}{[1][29][]29}}
\newlabel{eje:5-limites}{{5.2}{29}{}{ejercicio.5.2}{}}
\newlabel{eje:5-limites@cref}{{[ejercicio][2][5]5.2}{[1][29][]29}}
\newlabel{eje:5-geometricas}{{5.3}{29}{}{ejercicio.5.3}{}}
\newlabel{eje:5-geometricas@cref}{{[ejercicio][3][5]5.3}{[1][29][]29}}
\newlabel{eje:5-pelota}{{5.4}{29}{}{ejercicio.5.4}{}}
\newlabel{eje:5-pelota@cref}{{[ejercicio][4][5]5.4}{[1][29][]29}}
\newlabel{eje:5-comparacion}{{5.5}{29}{}{ejercicio.5.5}{}}
\newlabel{eje:5-comparacion@cref}{{[ejercicio][5][5]5.5}{[1][29][]29}}
\newlabel{eje:5-cociente}{{5.6}{29}{}{ejercicio.5.6}{}}
\newlabel{eje:5-cociente@cref}{{[ejercicio][6][5]5.6}{[1][29][]29}}
\newlabel{eje:5-integral}{{5.7}{30}{}{ejercicio.5.7}{}}
\newlabel{eje:5-integral@cref}{{[ejercicio][7][5]5.7}{[1][30][]30}}
\newlabel{eje:5-armonica}{{5.8}{30}{}{ejercicio.5.8}{}}
\newlabel{eje:5-armonica@cref}{{[ejercicio][8][5]5.8}{[1][30][]30}}
\newlabel{eje:5-alternadas}{{5.9}{30}{}{ejercicio.5.9}{}}
\newlabel{eje:5-alternadas@cref}{{[ejercicio][9][5]5.9}{[1][30][]30}}
\newlabel{eje:5-error-alternada}{{5.10}{30}{}{ejercicio.5.10}{}}
\newlabel{eje:5-error-alternada@cref}{{[ejercicio][10][5]5.10}{[1][30][]30}}
\newlabel{eje:5-descuento}{{5.11}{30}{}{ejercicio.5.11}{}}
\newlabel{eje:5-descuento@cref}{{[ejercicio][11][5]5.11}{[1][30][]30}}
\newlabel{eje:5-armonica-dem}{{5.12}{30}{}{ejercicio.5.12}{}}
\newlabel{eje:5-armonica-dem@cref}{{[ejercicio][12][5]5.12}{[1][30][]30}}
\newlabel{eje:5-cociente-dem}{{5.13}{30}{}{ejercicio.5.13}{}}
\newlabel{eje:5-cociente-dem@cref}{{[ejercicio][13][5]5.13}{[1][30][]30}}
\newlabel{eje:5-resto}{{5.14}{30}{}{ejercicio.5.14}{}}
\newlabel{eje:5-resto@cref}{{[ejercicio][14][5]5.14}{[1][30][]30}}
\newlabel{eje:5-euler}{{5.15}{30}{}{ejercicio.5.15}{}}
\newlabel{eje:5-euler@cref}{{[ejercicio][15][5]5.15}{[1][30][]30}}
\newlabel{eje:5-heron}{{5.16}{30}{}{ejercicio.5.16}{}}
\newlabel{eje:5-heron@cref}{{[ejercicio][16][5]5.16}{[1][30][]30}}
\newlabel{ej:radios}{{6.4}{32}{}{ejemplo.6.4}{}}
\newlabel{ej:radios@cref}{{[ejemplo][4][]6.4}{[1][32][]32}}
\newlabel{ej:desde-geometrica}{{6.7}{32}{Series obtenidas de la geométrica}{ejemplo.6.7}{}}
\newlabel{ej:desde-geometrica@cref}{{[ejemplo][7][]6.7}{[1][32][]32}}
\newlabel{obs:6-errores}{{6.16}{35}{Errores frecuentes}{observacion.6.16}{}}
\newlabel{obs:6-errores@cref}{{[observacion][16][]6.16}{[1][35][]35}}
\newlabel{eje:6-vf}{{6.1}{35}{}{ejercicio.6.1}{}}
\newlabel{eje:6-vf@cref}{{[ejercicio][1][6]6.1}{[1][35][]35}}
\newlabel{eje:6-radio}{{6.2}{35}{}{ejercicio.6.2}{}}
\newlabel{eje:6-radio@cref}{{[ejercicio][2][6]6.2}{[1][35][]35}}
\newlabel{eje:6-e2x}{{6.3}{35}{}{ejercicio.6.3}{}}
\newlabel{eje:6-e2x@cref}{{[ejercicio][3][6]6.3}{[1][35][]35}}
\newlabel{eje:6-taylor-basicas}{{6.4}{36}{}{ejercicio.6.4}{}}
\newlabel{eje:6-taylor-basicas@cref}{{[ejercicio][4][6]6.4}{[1][36][]36}}
\newlabel{eje:6-derivar-geom}{{6.5}{36}{}{ejercicio.6.5}{}}
\newlabel{eje:6-derivar-geom@cref}{{[ejercicio][5][6]6.5}{[1][36][]36}}
\newlabel{eje:6-productos}{{6.6}{36}{}{ejercicio.6.6}{}}
\newlabel{eje:6-productos@cref}{{[ejercicio][6][6]6.6}{[1][36][]36}}
\newlabel{eje:6-limites}{{6.7}{36}{}{ejercicio.6.7}{}}
\newlabel{eje:6-limites@cref}{{[ejercicio][7][6]6.7}{[1][36][]36}}
\newlabel{eje:6-cos1}{{6.8}{36}{}{ejercicio.6.8}{}}
\newlabel{eje:6-cos1@cref}{{[ejercicio][8][6]6.8}{[1][36][]36}}
\newlabel{eje:6-aproximar}{{6.9}{36}{}{ejercicio.6.9}{}}
\newlabel{eje:6-aproximar@cref}{{[ejercicio][9][6]6.9}{[1][36][]36}}
\newlabel{eje:6-seno-dem}{{6.10}{36}{}{ejercicio.6.10}{}}
\newlabel{eje:6-seno-dem@cref}{{[ejercicio][10][6]6.10}{[1][36][]36}}
\newlabel{eje:6-arctan}{{6.11}{36}{}{ejercicio.6.11}{}}
\newlabel{eje:6-arctan@cref}{{[ejercicio][11][6]6.11}{[1][36][]36}}
\newlabel{eje:6-binomial-dem}{{6.12}{36}{}{ejercicio.6.12}{}}
\newlabel{eje:6-binomial-dem@cref}{{[ejercicio][12][6]6.12}{[1][36][]36}}
\newlabel{eje:6-limite-exp}{{6.13}{36}{}{ejercicio.6.13}{}}
\newlabel{eje:6-limite-exp@cref}{{[ejercicio][13][6]6.13}{[1][36][]36}}
\newlabel{eje:6-no-analitica}{{6.14}{37}{}{ejercicio.6.14}{}}
\newlabel{eje:6-no-analitica@cref}{{[ejercicio][14][6]6.14}{[1][37][]37}}
\newlabel{ap:soluciones}{{A}{38}{Soluciones de los ejercicios}{appendix.A}{}}
\newlabel{ap:soluciones@cref}{{[appendix][1][2147483647]A}{[1][38][]38}}
\newlabel{sol:eje:1-inmediatas}{{A.1}{38}{Tema 1}{section*.3}{}}
\newlabel{sol:eje:1-inmediatas@cref}{{[subappendix][1][2147483647,1]A.1}{[1][38][]38}}
\newlabel{sol:eje:1-vf}{{A.1}{38}{Tema 1}{section*.4}{}}
\newlabel{sol:eje:1-vf@cref}{{[subappendix][1][2147483647,1]A.1}{[1][38][]38}}
\newlabel{sol:eje:1-sumas-x3}{{A.1}{39}{Tema 1}{section*.5}{}}
\newlabel{sol:eje:1-sumas-x3@cref}{{[subappendix][1][2147483647,1]A.1}{[1][39][]39}}
\newlabel{sol:eje:1-sumas-strang}{{A.1}{39}{Tema 1}{section*.6}{}}
\newlabel{sol:eje:1-sumas-strang@cref}{{[subappendix][1][2147483647,1]A.1}{[1][39][]39}}
\newlabel{sol:eje:1-limite-x3}{{A.1}{39}{Tema 1}{section*.7}{}}
\newlabel{sol:eje:1-limite-x3@cref}{{[subappendix][1][2147483647,1]A.1}{[1][39][]39}}
\newlabel{sol:eje:1-limite-seno}{{A.1}{40}{Tema 1}{section*.8}{}}
\newlabel{sol:eje:1-limite-seno@cref}{{[subappendix][1][2147483647,1]A.1}{[1][40][]40}}
\newlabel{sol:eje:1-limites-sumas}{{A.1}{40}{Tema 1}{section*.9}{}}
\newlabel{sol:eje:1-limites-sumas@cref}{{[subappendix][1][2147483647,1]A.1}{[1][40][]40}}
\newlabel{sol:eje:1-condicion}{{A.1}{40}{Tema 1}{section*.10}{}}
\newlabel{sol:eje:1-condicion@cref}{{[subappendix][1][2147483647,1]A.1}{[1][40][]40}}
\newlabel{sol:eje:1-estimar}{{A.1}{41}{Tema 1}{section*.11}{}}
\newlabel{sol:eje:1-estimar@cref}{{[subappendix][1][2147483647,1]A.1}{[1][40][]41}}
\newlabel{sol:eje:1-exp-sumas}{{A.1}{41}{Tema 1}{section*.12}{}}
\newlabel{sol:eje:1-exp-sumas@cref}{{[subappendix][1][2147483647,1]A.1}{[1][41][]41}}
\newlabel{sol:eje:1-diferencia}{{A.1}{41}{Tema 1}{section*.13}{}}
\newlabel{sol:eje:1-diferencia@cref}{{[subappendix][1][2147483647,1]A.1}{[1][41][]41}}
\newlabel{sol:eje:1-positiva}{{A.1}{42}{Tema 1}{section*.14}{}}
\newlabel{sol:eje:1-positiva@cref}{{[subappendix][1][2147483647,1]A.1}{[1][42][]42}}
\newlabel{sol:eje:1-un-punto}{{A.1}{42}{Tema 1}{section*.15}{}}
\newlabel{sol:eje:1-un-punto@cref}{{[subappendix][1][2147483647,1]A.1}{[1][42][]42}}
\newlabel{sol:eje:1-nula}{{A.1}{43}{Tema 1}{section*.16}{}}
\newlabel{sol:eje:1-nula@cref}{{[subappendix][1][2147483647,1]A.1}{[1][42][]43}}
\newlabel{sol:eje:1-racionales-irr}{{A.1}{43}{Tema 1}{section*.17}{}}
\newlabel{sol:eje:1-racionales-irr@cref}{{[subappendix][1][2147483647,1]A.1}{[1][43][]43}}
\newlabel{sol:eje:2-conceptos}{{A.2}{44}{Tema 2}{section*.18}{}}
\newlabel{sol:eje:2-conceptos@cref}{{[subappendix][2][2147483647,1]A.2}{[1][43][]44}}
\newlabel{sol:eje:2-error-barrow}{{A.2}{44}{Tema 2}{section*.19}{}}
\newlabel{sol:eje:2-error-barrow@cref}{{[subappendix][2][2147483647,1]A.2}{[1][44][]44}}
\newlabel{sol:eje:2-barrow}{{A.2}{44}{Tema 2}{section*.20}{}}
\newlabel{sol:eje:2-barrow@cref}{{[subappendix][2][2147483647,1]A.2}{[1][44][]44}}
\newlabel{sol:eje:2-arco-seno}{{A.2}{45}{Tema 2}{section*.21}{}}
\newlabel{sol:eje:2-arco-seno@cref}{{[subappendix][2][2147483647,1]A.2}{[1][45][]45}}
\newlabel{sol:eje:2-parabola-eje}{{A.2}{45}{Tema 2}{section*.22}{}}
\newlabel{sol:eje:2-parabola-eje@cref}{{[subappendix][2][2147483647,1]A.2}{[1][45][]45}}
\newlabel{sol:eje:2-derivar}{{A.2}{45}{Tema 2}{section*.23}{}}
\newlabel{sol:eje:2-derivar@cref}{{[subappendix][2][2147483647,1]A.2}{[1][45][]45}}
\newlabel{sol:eje:2-hallar-f}{{A.2}{46}{Tema 2}{section*.24}{}}
\newlabel{sol:eje:2-hallar-f@cref}{{[subappendix][2][2147483647,1]A.2}{[1][46][]46}}
\newlabel{sol:eje:2-valor-medio}{{A.2}{46}{Tema 2}{section*.25}{}}
\newlabel{sol:eje:2-valor-medio@cref}{{[subappendix][2][2147483647,1]A.2}{[1][46][]46}}
\newlabel{sol:eje:2-estudio-F}{{A.2}{46}{Tema 2}{section*.26}{}}
\newlabel{sol:eje:2-estudio-F@cref}{{[subappendix][2][2147483647,1]A.2}{[1][46][]46}}
\newlabel{sol:eje:2-limite}{{A.2}{47}{Tema 2}{section*.27}{}}
\newlabel{sol:eje:2-limite@cref}{{[subappendix][2][2147483647,1]A.2}{[1][46][]47}}
\newlabel{sol:eje:2-par-impar}{{A.2}{47}{Tema 2}{section*.28}{}}
\newlabel{sol:eje:2-par-impar@cref}{{[subappendix][2][2147483647,1]A.2}{[1][47][]47}}
\newlabel{sol:eje:2-logaritmo}{{A.2}{47}{Tema 2}{section*.29}{}}
\newlabel{sol:eje:2-logaritmo@cref}{{[subappendix][2][2147483647,1]A.2}{[1][47][]47}}
\newlabel{sol:eje:2-periodica}{{A.2}{48}{Tema 2}{section*.30}{}}
\newlabel{sol:eje:2-periodica@cref}{{[subappendix][2][2147483647,1]A.2}{[1][48][]48}}
\newlabel{sol:eje:2-simetrica-f}{{A.2}{48}{Tema 2}{section*.31}{}}
\newlabel{sol:eje:2-simetrica-f@cref}{{[subappendix][2][2147483647,1]A.2}{[1][48][]48}}
\newlabel{sol:eje:2-tvm-ponderado}{{A.2}{48}{Tema 2}{section*.32}{}}
\newlabel{sol:eje:2-tvm-ponderado@cref}{{[subappendix][2][2147483647,1]A.2}{[1][48][]48}}
\newlabel{sol:eje:3-elegir}{{A.3}{49}{Tema 3}{section*.33}{}}
\newlabel{sol:eje:3-elegir@cref}{{[subappendix][3][2147483647,1]A.3}{[1][49][]49}}
\newlabel{sol:eje:3-sustitucion}{{A.3}{49}{Tema 3}{section*.34}{}}
\newlabel{sol:eje:3-sustitucion@cref}{{[subappendix][3][2147483647,1]A.3}{[1][49][]49}}
\newlabel{sol:eje:3-limites}{{A.3}{50}{Tema 3}{section*.35}{}}
\newlabel{sol:eje:3-limites@cref}{{[subappendix][3][2147483647,1]A.3}{[1][49][]50}}
\newlabel{sol:eje:3-partes}{{A.3}{50}{Tema 3}{section*.36}{}}
\newlabel{sol:eje:3-partes@cref}{{[subappendix][3][2147483647,1]A.3}{[1][50][]50}}
\newlabel{sol:eje:3-ciclicas}{{A.3}{50}{Tema 3}{section*.37}{}}
\newlabel{sol:eje:3-ciclicas@cref}{{[subappendix][3][2147483647,1]A.3}{[1][50][]50}}
\newlabel{sol:eje:3-fracciones}{{A.3}{51}{Tema 3}{section*.38}{}}
\newlabel{sol:eje:3-fracciones@cref}{{[subappendix][3][2147483647,1]A.3}{[1][50][]51}}
\newlabel{sol:eje:3-fracciones-cuad}{{A.3}{51}{Tema 3}{section*.39}{}}
\newlabel{sol:eje:3-fracciones-cuad@cref}{{[subappendix][3][2147483647,1]A.3}{[1][51][]51}}
\newlabel{sol:eje:3-racional-impropia}{{A.3}{51}{Tema 3}{section*.40}{}}
\newlabel{sol:eje:3-racional-impropia@cref}{{[subappendix][3][2147483647,1]A.3}{[1][51][]51}}
\newlabel{sol:eje:3-trig}{{A.3}{52}{Tema 3}{section*.41}{}}
\newlabel{sol:eje:3-trig@cref}{{[subappendix][3][2147483647,1]A.3}{[1][51][]52}}
\newlabel{sol:eje:3-trig-sust}{{A.3}{52}{Tema 3}{section*.42}{}}
\newlabel{sol:eje:3-trig-sust@cref}{{[subappendix][3][2147483647,1]A.3}{[1][52][]52}}
\newlabel{sol:eje:3-curso}{{A.3}{52}{Tema 3}{section*.43}{}}
\newlabel{sol:eje:3-curso@cref}{{[subappendix][3][2147483647,1]A.3}{[1][52][]52}}
\newlabel{sol:eje:3-reduccion}{{A.3}{53}{Tema 3}{section*.44}{}}
\newlabel{sol:eje:3-reduccion@cref}{{[subappendix][3][2147483647,1]A.3}{[1][52][]53}}
\newlabel{sol:eje:3-wallis}{{A.3}{53}{Tema 3}{section*.45}{}}
\newlabel{sol:eje:3-wallis@cref}{{[subappendix][3][2147483647,1]A.3}{[1][53][]53}}
\newlabel{sol:eje:3-beta}{{A.3}{53}{Tema 3}{section*.46}{}}
\newlabel{sol:eje:3-beta@cref}{{[subappendix][3][2147483647,1]A.3}{[1][53][]53}}
\newlabel{sol:eje:3-sustitucion-falsa}{{A.3}{54}{Tema 3}{section*.47}{}}
\newlabel{sol:eje:3-sustitucion-falsa@cref}{{[subappendix][3][2147483647,1]A.3}{[1][54][]54}}
\newlabel{sol:eje:4-clasificar}{{A.4}{54}{Tema 4}{section*.48}{}}
\newlabel{sol:eje:4-clasificar@cref}{{[subappendix][4][2147483647,1]A.4}{[1][54][]54}}
\newlabel{sol:eje:4-calcular}{{A.4}{55}{Tema 4}{section*.49}{}}
\newlabel{sol:eje:4-calcular@cref}{{[subappendix][4][2147483647,1]A.4}{[1][55][]55}}
\newlabel{sol:eje:4-logaritmos}{{A.4}{55}{Tema 4}{section*.50}{}}
\newlabel{sol:eje:4-logaritmos@cref}{{[subappendix][4][2147483647,1]A.4}{[1][55][]55}}
\newlabel{sol:eje:4-comparar}{{A.4}{55}{Tema 4}{section*.51}{}}
\newlabel{sol:eje:4-comparar@cref}{{[subappendix][4][2147483647,1]A.4}{[1][55][]55}}
\newlabel{sol:eje:4-seno}{{A.4}{56}{Tema 4}{section*.52}{}}
\newlabel{sol:eje:4-seno@cref}{{[subappendix][4][2147483647,1]A.4}{[1][56][]56}}
\newlabel{sol:eje:4-gamma-valores}{{A.4}{56}{Tema 4}{section*.53}{}}
\newlabel{sol:eje:4-gamma-valores@cref}{{[subappendix][4][2147483647,1]A.4}{[1][56][]56}}
\newlabel{sol:eje:4-gamma-medio}{{A.4}{56}{Tema 4}{section*.54}{}}
\newlabel{sol:eje:4-gamma-medio@cref}{{[subappendix][4][2147483647,1]A.4}{[1][56][]56}}
\newlabel{sol:eje:4-erf}{{A.4}{57}{Tema 4}{section*.55}{}}
\newlabel{sol:eje:4-erf@cref}{{[subappendix][4][2147483647,1]A.4}{[1][57][]57}}
\newlabel{sol:eje:4-normal}{{A.4}{57}{Tema 4}{section*.56}{}}
\newlabel{sol:eje:4-normal@cref}{{[subappendix][4][2147483647,1]A.4}{[1][57][]57}}
\newlabel{sol:eje:4-gamma-conv}{{A.4}{57}{Tema 4}{section*.57}{}}
\newlabel{sol:eje:4-gamma-conv@cref}{{[subappendix][4][2147483647,1]A.4}{[1][57][]57}}
\newlabel{sol:eje:4-limite-comp}{{A.4}{58}{Tema 4}{section*.58}{}}
\newlabel{sol:eje:4-limite-comp@cref}{{[subappendix][4][2147483647,1]A.4}{[1][58][]58}}
\newlabel{sol:eje:4-dirichlet}{{A.4}{58}{Tema 4}{section*.59}{}}
\newlabel{sol:eje:4-dirichlet@cref}{{[subappendix][4][2147483647,1]A.4}{[1][58][]58}}
\newlabel{sol:eje:4-potencias}{{A.4}{59}{Tema 4}{section*.60}{}}
\newlabel{sol:eje:4-potencias@cref}{{[subappendix][4][2147483647,1]A.4}{[1][59][]59}}
\newlabel{sol:eje:4-decreciente}{{A.4}{59}{Tema 4}{section*.61}{}}
\newlabel{sol:eje:4-decreciente@cref}{{[subappendix][4][2147483647,1]A.4}{[1][59][]59}}
\newlabel{sol:eje:4-wallis-gauss}{{A.4}{59}{Tema 4}{section*.62}{}}
\newlabel{sol:eje:4-wallis-gauss@cref}{{[subappendix][4][2147483647,1]A.4}{[1][59][]59}}
\newlabel{sol:eje:5-vf}{{A.5}{60}{Tema 5}{section*.63}{}}
\newlabel{sol:eje:5-vf@cref}{{[subappendix][5][2147483647,1]A.5}{[1][60][]60}}
\newlabel{sol:eje:5-limites}{{A.5}{60}{Tema 5}{section*.64}{}}
\newlabel{sol:eje:5-limites@cref}{{[subappendix][5][2147483647,1]A.5}{[1][60][]60}}
\newlabel{sol:eje:5-geometricas}{{A.5}{61}{Tema 5}{section*.65}{}}
\newlabel{sol:eje:5-geometricas@cref}{{[subappendix][5][2147483647,1]A.5}{[1][61][]61}}
\newlabel{sol:eje:5-pelota}{{A.5}{61}{Tema 5}{section*.66}{}}
\newlabel{sol:eje:5-pelota@cref}{{[subappendix][5][2147483647,1]A.5}{[1][61][]61}}
\newlabel{sol:eje:5-comparacion}{{A.5}{61}{Tema 5}{section*.67}{}}
\newlabel{sol:eje:5-comparacion@cref}{{[subappendix][5][2147483647,1]A.5}{[1][61][]61}}
\newlabel{sol:eje:5-cociente}{{A.5}{62}{Tema 5}{section*.68}{}}
\newlabel{sol:eje:5-cociente@cref}{{[subappendix][5][2147483647,1]A.5}{[1][62][]62}}
\newlabel{sol:eje:5-integral}{{A.5}{62}{Tema 5}{section*.69}{}}
\newlabel{sol:eje:5-integral@cref}{{[subappendix][5][2147483647,1]A.5}{[1][62][]62}}
\newlabel{sol:eje:5-armonica}{{A.5}{62}{Tema 5}{section*.70}{}}
\newlabel{sol:eje:5-armonica@cref}{{[subappendix][5][2147483647,1]A.5}{[1][62][]62}}
\newlabel{sol:eje:5-alternadas}{{A.5}{63}{Tema 5}{section*.71}{}}
\newlabel{sol:eje:5-alternadas@cref}{{[subappendix][5][2147483647,1]A.5}{[1][63][]63}}
\newlabel{sol:eje:5-error-alternada}{{A.5}{63}{Tema 5}{section*.72}{}}
\newlabel{sol:eje:5-error-alternada@cref}{{[subappendix][5][2147483647,1]A.5}{[1][63][]63}}
\newlabel{sol:eje:5-descuento}{{A.5}{63}{Tema 5}{section*.73}{}}
\newlabel{sol:eje:5-descuento@cref}{{[subappendix][5][2147483647,1]A.5}{[1][63][]63}}
\newlabel{sol:eje:5-armonica-dem}{{A.5}{64}{Tema 5}{section*.74}{}}
\newlabel{sol:eje:5-armonica-dem@cref}{{[subappendix][5][2147483647,1]A.5}{[1][64][]64}}
\newlabel{sol:eje:5-cociente-dem}{{A.5}{64}{Tema 5}{section*.75}{}}
\newlabel{sol:eje:5-cociente-dem@cref}{{[subappendix][5][2147483647,1]A.5}{[1][64][]64}}
\newlabel{sol:eje:5-resto}{{A.5}{64}{Tema 5}{section*.76}{}}
\newlabel{sol:eje:5-resto@cref}{{[subappendix][5][2147483647,1]A.5}{[1][64][]64}}
\newlabel{sol:eje:5-euler}{{A.5}{65}{Tema 5}{section*.77}{}}
\newlabel{sol:eje:5-euler@cref}{{[subappendix][5][2147483647,1]A.5}{[1][65][]65}}
\newlabel{sol:eje:5-heron}{{A.5}{65}{Tema 5}{section*.78}{}}
\newlabel{sol:eje:5-heron@cref}{{[subappendix][5][2147483647,1]A.5}{[1][65][]65}}
\newlabel{sol:eje:6-vf}{{A.6}{66}{Tema 6}{section*.79}{}}
\newlabel{sol:eje:6-vf@cref}{{[subappendix][6][2147483647,1]A.6}{[1][66][]66}}
\newlabel{sol:eje:6-radio}{{A.6}{66}{Tema 6}{section*.80}{}}
\newlabel{sol:eje:6-radio@cref}{{[subappendix][6][2147483647,1]A.6}{[1][66][]66}}
\newlabel{sol:eje:6-e2x}{{A.6}{66}{Tema 6}{section*.81}{}}
\newlabel{sol:eje:6-e2x@cref}{{[subappendix][6][2147483647,1]A.6}{[1][66][]66}}
\newlabel{sol:eje:6-taylor-basicas}{{A.6}{67}{Tema 6}{section*.82}{}}
\newlabel{sol:eje:6-taylor-basicas@cref}{{[subappendix][6][2147483647,1]A.6}{[1][67][]67}}
\newlabel{sol:eje:6-derivar-geom}{{A.6}{67}{Tema 6}{section*.83}{}}
\newlabel{sol:eje:6-derivar-geom@cref}{{[subappendix][6][2147483647,1]A.6}{[1][67][]67}}
\newlabel{sol:eje:6-productos}{{A.6}{67}{Tema 6}{section*.84}{}}
\newlabel{sol:eje:6-productos@cref}{{[subappendix][6][2147483647,1]A.6}{[1][67][]67}}
\newlabel{sol:eje:6-limites}{{A.6}{68}{Tema 6}{section*.85}{}}
\newlabel{sol:eje:6-limites@cref}{{[subappendix][6][2147483647,1]A.6}{[1][68][]68}}
\newlabel{sol:eje:6-cos1}{{A.6}{68}{Tema 6}{section*.86}{}}
\newlabel{sol:eje:6-cos1@cref}{{[subappendix][6][2147483647,1]A.6}{[1][68][]68}}
\newlabel{sol:eje:6-aproximar}{{A.6}{68}{Tema 6}{section*.87}{}}
\newlabel{sol:eje:6-aproximar@cref}{{[subappendix][6][2147483647,1]A.6}{[1][68][]68}}
\newlabel{sol:eje:6-seno-dem}{{A.6}{69}{Tema 6}{section*.88}{}}
\newlabel{sol:eje:6-seno-dem@cref}{{[subappendix][6][2147483647,1]A.6}{[1][69][]69}}
\newlabel{sol:eje:6-arctan}{{A.6}{69}{Tema 6}{section*.89}{}}
\newlabel{sol:eje:6-arctan@cref}{{[subappendix][6][2147483647,1]A.6}{[1][69][]69}}
\newlabel{sol:eje:6-binomial-dem}{{A.6}{70}{Tema 6}{section*.90}{}}
\newlabel{sol:eje:6-binomial-dem@cref}{{[subappendix][6][2147483647,1]A.6}{[1][70][]70}}
\newlabel{sol:eje:6-limite-exp}{{A.6}{70}{Tema 6}{section*.91}{}}
\newlabel{sol:eje:6-limite-exp@cref}{{[subappendix][6][2147483647,1]A.6}{[1][70][]70}}
\newlabel{sol:eje:6-no-analitica}{{A.6}{71}{Tema 6}{section*.92}{}}
\newlabel{sol:eje:6-no-analitica@cref}{{[subappendix][6][2147483647,1]A.6}{[1][70][]71}}
\makeatother

\begin{document}
\frontmatter
\begin{titlepage}
\centering
\vspace*{2.5cm}
{\large\scshape Anatomía de un LLM}\par
\vspace{0.4cm}
{\normalsize Fundamentos matemáticos de los modelos de lenguaje}\par
\vspace{2.2cm}
\rule{\textwidth}{0.8pt}\par\vspace{0.5cm}
{\small\scshape Fase 0 \textperiodcentered\ Preliminares matemáticos}\par\vspace{0.3cm}
{\Huge\bfseries Módulo 0.6\par}\vspace{0.3cm}
{\LARGE Cálculo integral y series\par}
\vspace{0.5cm}\rule{\textwidth}{0.4pt}\par
\vspace{1.6cm}
{\large Resumen del módulo}\par
\vfill
{\small Edición de octubre de 2026}\par
\end{titlepage}

\chapter*{Sobre este resumen}
\addcontentsline{toc}{chapter}{Sobre este resumen}
\markboth{Sobre este resumen}{}

Este documento reúne lo indispensable de los apuntes del Módulo 0.6: el texto expositivo, todas las definiciones y todos los resultados (teoremas, proposiciones, lemas y corolarios), las tablas y las observaciones de contenido. Se han suprimido las demostraciones, los ejemplos, las lecturas, las figuras, los ejercicios y sus soluciones. La numeración de temas, secciones, enunciados y ecuaciones es la misma que en los apuntes, de modo que cualquier referencia puede buscarse allí; las que remiten a elementos suprimidos (un ejemplo, una figura, una ecuación de una demostración) conservan su número en los apuntes.

\section*{Notación}
\begin{center}
\begin{tabular}{@{}ll@{}}
\toprule
Símbolo & Significado \\
\midrule
$\N=\{1,2,3,\dots\}$, $\N_0=\N\cup\{0\}$ & números naturales, sin y con el cero \\
$P=\{x_0,\dots,x_n\}$, $\Delta x_i=x_i-x_{i-1}$ & partición de $[a,b]$ y longitud de sus subintervalos \\
$\Lsum(f,P)$, $\Usum(f,P)$ & suma inferior y suma superior de $f$ para $P$ \\
$\Rsum(f,P,\xi)$ & suma de Riemann con puntos intermedios $\xi_i$ \\
$\int f(x)\dif x$ & primitiva (integral indefinida), definida salvo constante \\
$\int_a^b f(x)\dif x$ & integral definida de $f$ en $[a,b]$ \\
$\eval{F(x)}{a}{b}$ & $F(b)-F(a)$ \\
$\ln$ & logaritmo natural \\
$\Gamma(s)$ & función Gamma, $\int_0^\infty x^{s-1}e^{-x}\dif x$ \\
$(a_n)_{n\ge1}$, $\sum a_n$ & sucesión y serie \\
$S_n=\sum_{k=1}^{n}a_k$ & suma parcial $n$-ésima \\
$H_n=\sum_{k=1}^{n}1/k$ & número armónico \\
$\binom{a}{n}=\frac{a(a-1)\cdots(a-n+1)}{n!}$ & coeficiente binomial generalizado, $a\in\R$ \\
$T_n f$, $R_n f$ & polinomio de Taylor de grado $n$ y su resto \\
\bottomrule
\end{tabular}
\end{center}


\tableofcontents

\mainmatter
% =====================================================================
\chapter{Primitivas e integral de Riemann}
\label{tema:1}
% =====================================================================

El cálculo integral responde a dos preguntas que, en apariencia, no tienen relación: qué función tiene una derivada dada y cuánto vale el área bajo una curva. Este tema plantea las dos por separado. La primera conduce a las primitivas; la segunda, a la integral de Riemann, definida como el número que queda atrapado entre las sumas inferiores y superiores de rectángulos. El teorema fundamental del \cref{tema:2} mostrará que ambas preguntas son la misma. La exposición sigue los capítulos~13 de Spivak~\cite{spivak} y~1 de Apostol~\cite{apostol}, con los ejemplos y ejercicios de las §5.1 y~5.2 de~\cite{openstax1}, de la §5.5 de Strang~\cite{strang} y de la parte~A de la unidad~3 del curso 18.01SC del MIT~\cite{mit1801}.

\section{Primitivas}

\begin{definicion}[Primitiva]\label{def:primitiva}
Sea $I\subseteq\R$ un intervalo y sea $f\colon I\to\R$. Una función $F\colon I\to\R$ es una \emph{primitiva} de $f$ en $I$ si $F$ es derivable en $I$ y $F'(x)=f(x)$ para todo $x\in I$. El conjunto de todas las primitivas de $f$ se denota $\int f(x)\dif x$ y se llama \emph{integral indefinida} de $f$.
\end{definicion}

\begin{proposicion}\label{prop:primitivas-constante}
Sean $F$ y $G$ dos primitivas de $f$ en un intervalo $I$. Existe una constante $C\in\R$ tal que $G(x)=F(x)+C$ para todo $x\in I$. En consecuencia, si $F$ es una primitiva, $\int f(x)\dif x=\{F+C: C\in\R\}$, lo que se escribe abreviadamente $\int f(x)\dif x=F(x)+C$.
\end{proposicion}

La hipótesis de que $I$ sea un intervalo es imprescindible. En $\R\setminus\{0\}$, la función $H(x)=0$ si $x<0$ y $H(x)=1$ si $x>0$ tiene derivada nula y no es constante; por eso las primitivas de $1/x$ en $\R\setminus\{0\}$ son $\ln\abs{x}+C_1$ para $x<0$ y $\ln\abs{x}+C_2$ para $x>0$, con constantes independientes.

Como la derivada es lineal, también lo es la integral indefinida: si $F$ y $G$ son primitivas de $f$ y $g$, entonces $\alpha F+\beta G$ es primitiva de $\alpha f+\beta g$. Leyendo al revés la tabla de derivadas del módulo~0.5 se obtienen las \emph{primitivas inmediatas} de la \cref{tab:inmediatas}; cada fila se comprueba derivando la columna derecha.

\begin{table}[h]
\centering
\small
\begin{tabular}{@{}lll@{}}
\toprule
$f(x)$ & $\int f(x)\dif x$ & Validez \\
\midrule
$x^{r}$, $r\neq-1$ & $\dfrac{x^{r+1}}{r+1}+C$ & $x>0$ (todo $x$ si $r\in\N_0$) \\[4pt]
$\dfrac1x$ & $\ln\abs{x}+C$ & cada intervalo de $\R\setminus\{0\}$ \\[4pt]
$e^{ax}$, $a\neq0$ & $\dfrac{e^{ax}}{a}+C$ & $\R$ \\[4pt]
$\sin(ax)$, $\cos(ax)$, $a\neq0$ & $-\dfrac{\cos(ax)}{a}+C$, \ $\dfrac{\sin(ax)}{a}+C$ & $\R$ \\[4pt]
$\dfrac{1}{\cos^{2}x}=1+\tan^{2}x$ & $\tan x+C$ & cada intervalo donde $\cos x\neq0$ \\[4pt]
$\dfrac{1}{1+x^{2}}$ & $\arctan x+C$ & $\R$ \\[4pt]
$\dfrac{1}{\sqrt{1-x^{2}}}$ & $\arcsin x+C$ & $(-1,1)$ \\[4pt]
$f(ax+b)$, $a\neq0$ & $\dfrac{1}{a}F(ax+b)+C$, con $F'=f$ & donde esté definida \\
\bottomrule
\end{tabular}
\caption{Primitivas inmediatas. La última fila es la regla de la cadena leída al revés.}
\label{tab:inmediatas}
\end{table}

\addtocounter{definicion}{1}%

\section{Particiones y sumas de Darboux}

La idea de la integral es aproximar el área bajo la gráfica de $f$ por rectángulos y refinar la aproximación. Para que la definición sea rigurosa conviene encerrar el área entre una aproximación por defecto y otra por exceso.

\begin{definicion}[Partición]\label{def:particion}
Sea $[a,b]$ un intervalo con $a<b$. Una \emph{partición} de $[a,b]$ es un conjunto finito $P=\{x_0,x_1,\dots,x_n\}$ con $a=x_0<x_1<\dots<x_n=b$. Se escribe $\Delta x_i=x_i-x_{i-1}$. La \emph{norma} de $P$ es $\norm{P}=\max_i\Delta x_i$. Una partición $Q$ es \emph{más fina} que $P$ si $P\subseteq Q$. La \emph{partición uniforme} de $n$ subintervalos es $x_i=a+i\,\frac{b-a}{n}$.
\end{definicion}

\begin{definicion}[Sumas inferior y superior]\label{def:sumas}
Sea $f\colon[a,b]\to\R$ acotada y sea $P$ una partición. Para cada $i$ se definen
$m_i=\inf\{f(x):x\in[x_{i-1},x_i]\}$ y $M_i=\sup\{f(x):x\in[x_{i-1},x_i]\}$, y
\[
\Lsum(f,P)=\sum_{i=1}^{n}m_i\,\Delta x_i,\qquad \Usum(f,P)=\sum_{i=1}^{n}M_i\,\Delta x_i .
\]
\end{definicion}

Si $f\ge0$, $\Lsum(f,P)$ es el área de una unión de rectángulos contenida en la región bajo la gráfica, y $\Usum(f,P)$ el de una unión que la contiene (\cref{fig:darboux}). Siempre $\Lsum(f,P)\le\Usum(f,P)$, porque $m_i\le M_i$.

\addtocounter{figure}{1}%

\begin{lema}[Refinamiento]\label{lem:refinamiento}
Si $Q$ es más fina que $P$, entonces $\Lsum(f,P)\le\Lsum(f,Q)$ y $\Usum(f,Q)\le\Usum(f,P)$.
\end{lema}

\begin{corolario}\label{cor:inferior-superior}
Para dos particiones cualesquiera $P$ y $Q$ de $[a,b]$, $\Lsum(f,P)\le\Usum(f,Q)$.
\end{corolario}

\section{La integral de Riemann}

Por el \cref{cor:inferior-superior}, el conjunto de las sumas inferiores está acotado superiormente por cualquier suma superior. Existen, por tanto, la \emph{integral inferior} $\underline{I}=\sup_P\Lsum(f,P)$ y la \emph{integral superior} $\overline{I}=\inf_P\Usum(f,P)$, y se cumple $\underline{I}\le\overline{I}$.

\begin{definicion}[Integral de Riemann]\label{def:integral}
Una función acotada $f\colon[a,b]\to\R$ es \emph{integrable} (en el sentido de Riemann) si $\underline{I}=\overline{I}$. Ese valor común es la \emph{integral definida} de $f$ en $[a,b]$ y se denota $\int_a^b f(x)\dif x$. Se conviene además que $\int_a^a f=0$ y $\int_b^a f=-\int_a^b f$.
\end{definicion}

La variable $x$ de $\int_a^b f(x)\dif x$ es muda: $\int_a^b f(x)\dif x=\int_a^b f(t)\dif t$. La integral es un número; la integral indefinida, un conjunto de funciones. El teorema fundamental del \cref{tema:2} relacionará ambas.

\begin{teorema}[Criterio de integrabilidad de Riemann]\label{teo:criterio}
Una función acotada $f\colon[a,b]\to\R$ es integrable si y solo si para cada $\varepsilon>0$ existe una partición $P$ con $\Usum(f,P)-\Lsum(f,P)<\varepsilon$.
\end{teorema}

\addtocounter{definicion}{1}%

\addtocounter{definicion}{1}%

\begin{definicion}[Suma de Riemann]\label{def:suma-riemann}
Sea $P$ una partición de $[a,b]$ y sea $\xi=(\xi_1,\dots,\xi_n)$ con $\xi_i\in[x_{i-1},x_i]$. La \emph{suma de Riemann} es $\Rsum(f,P,\xi)=\sum_{i=1}^{n}f(\xi_i)\,\Delta x_i$. Con $\xi_i=x_{i-1}$, $\xi_i=x_i$ o $\xi_i=\frac{x_{i-1}+x_i}{2}$ se obtienen las sumas por la izquierda, por la derecha y del punto medio.
\end{definicion}

\begin{proposicion}\label{prop:sumas-riemann}
Sea $f$ integrable en $[a,b]$ y sea $(P_n)$ una sucesión de particiones con $\Usum(f,P_n)-\Lsum(f,P_n)\to0$. Para cualquier elección de puntos intermedios $\xi^{(n)}$, $\Rsum(f,P_n,\xi^{(n)})\to\int_a^b f$.
\end{proposicion}

La \cref{prop:sumas-riemann} permite calcular límites de sumas reconociéndolas como sumas de Riemann. Por ejemplo, con la partición uniforme de $[0,1]$ y $\xi_k=k/n$,
\[
\lim_{n\to\infty}\frac1n\sum_{k=1}^{n}\Bigl(\frac{k}{n}\Bigr)^{2}=\int_0^1x^{2}\dif x=\frac13 .
\]

\section{Funciones integrables}

\begin{teorema}[Las funciones monótonas son integrables]\label{teo:monotona}
Toda función monótona $f\colon[a,b]\to\R$ es integrable. Además, con la partición uniforme $P_n$, $\Usum(f,P_n)-\Lsum(f,P_n)=\frac{b-a}{n}\abs{f(b)-f(a)}$.
\end{teorema}

\begin{teorema}[Las funciones continuas son integrables]\label{teo:continua}
Toda función continua $f\colon[a,b]\to\R$ es integrable, y para la partición uniforme $\Usum(f,P_n)-\Lsum(f,P_n)\to0$.
\end{teorema}

Las mismas ideas prueban que una función acotada con un número finito de discontinuidades también es integrable, y que su integral no cambia si se modifica su valor en un número finito de puntos (\cref{eje:1-un-punto}).

\section{Propiedades de la integral}

\begin{teorema}[Propiedades de la integral]\label{teo:propiedades}
Sean $f,g$ integrables en $[a,b]$ y sean $\alpha,\beta\in\R$.
\begin{enumerate}
  \item (\emph{Linealidad}) $\alpha f+\beta g$ es integrable y $\int_a^b(\alpha f+\beta g)=\alpha\int_a^b f+\beta\int_a^b g$.
  \item (\emph{Aditividad}) Si $a<c<b$, $f$ es integrable en $[a,c]$ y en $[c,b]$, y $\int_a^b f=\int_a^c f+\int_c^b f$. Con el convenio de la \cref{def:integral}, la igualdad vale para cualesquiera $a,b,c$ en un intervalo donde $f$ sea integrable.
  \item (\emph{Monotonía}) Si $f\le g$ en $[a,b]$, entonces $\int_a^b f\le\int_a^b g$.
  \item (\emph{Acotación}) Si $m\le f(x)\le M$ en $[a,b]$, entonces $m(b-a)\le\int_a^b f\le M(b-a)$.
  \item (\emph{Valor absoluto}) $\abs{f}$ es integrable y $\bigl|\int_a^b f\bigr|\le\int_a^b\abs{f}$.
\end{enumerate}
\end{teorema}

\addtocounter{definicion}{1}%

\addtocounter{definicion}{1}%

\addtocounter{definicion}{1}%


% =====================================================================
\chapter{Teorema fundamental del cálculo, áreas y valor medio}
\label{tema:2}
% =====================================================================

El \cref{tema:1} definió la integral como un límite de sumas, una definición rigurosa pero poco práctica: calcular $\int_0^1x^{3}\dif x$ por sumas ya exige conocer $\sum k^{3}$. El teorema fundamental del cálculo cambia por completo la situación, porque reduce el cálculo de una integral al de una primitiva y, a la vez, garantiza que toda función continua tiene primitiva. Este tema lo demuestra en sus dos partes y lo aplica a áreas, valores medios y funciones definidas por integrales, como el logaritmo. Se siguen el capítulo~14 de Spivak~\cite{spivak}, el capítulo~5 de Apostol~\cite{apostol}, la §5.3 de~\cite{openstax1}, las §5.6 y~5.7 de Strang~\cite{strang} y las partes~A a~C de la unidad~3 del curso 18.01SC~\cite{mit1801}.

\section{El teorema del valor medio integral}

\begin{teorema}[Valor medio integral]\label{teo:tvm-integral}
Sea $f$ continua en $[a,b]$. Existe $c\in[a,b]$ tal que
\[
\int_a^b f(x)\dif x=f(c)\,(b-a).
\]
\end{teorema}

\begin{definicion}[Valor medio de una función]\label{def:valor-medio}
Sea $f$ integrable en $[a,b]$. Su \emph{valor medio} en $[a,b]$ es
$\displaystyle\bar f=\frac{1}{b-a}\int_a^b f(x)\dif x$.
\end{definicion}

Geométricamente, si $f\ge0$, el rectángulo de base $[a,b]$ y altura $\bar f$ tiene la misma área que la región bajo la gráfica. El \cref{teo:tvm-integral} dice que una función continua toma su valor medio en algún punto. El nombre procede de la analogía con la media de $n$ números: $\bar f$ es el límite de las medias $\frac1n\sum_{k=1}^{n}f(x_k)$ en puntos equiespaciados, que son sumas de Riemann divididas por $b-a$.

\section{El teorema fundamental del cálculo}

\begin{teorema}[Teorema fundamental del cálculo, primera parte]\label{teo:tfc1}
Sea $f$ integrable en $[a,b]$ y sea $F(x)=\int_a^x f(t)\dif t$ para $x\in[a,b]$.
\begin{enumerate}
  \item $F$ es continua en $[a,b]$; de hecho, si $\abs{f}\le K$, entonces $\abs{F(y)-F(x)}\le K\abs{y-x}$.
  \item Si $f$ es continua en un punto $c\in(a,b)$, entonces $F$ es derivable en $c$ y $F'(c)=f(c)$.
\end{enumerate}
\end{teorema}

\begin{corolario}\label{cor:existencia-primitiva}
Toda función continua en un intervalo $I$ tiene primitiva en $I$: para cualquier $a\in I$, $F(x)=\int_a^x f(t)\dif t$ lo es.
\end{corolario}

\begin{corolario}[Derivación de integrales con límites variables]\label{cor:leibniz}
Sea $f$ continua en un intervalo $I$ y sean $g,h$ derivables con valores en $I$. Entonces
\[
\frac{d}{dx}\int_{h(x)}^{g(x)}f(t)\dif t=f\bigl(g(x)\bigr)\,g'(x)-f\bigl(h(x)\bigr)\,h'(x).
\]
\end{corolario}

\begin{teorema}[Teorema fundamental del cálculo, segunda parte, o regla de Barrow]\label{teo:barrow}
Sea $f$ integrable en $[a,b]$ y sea $F$ continua en $[a,b]$, derivable en $(a,b)$, con $F'(x)=f(x)$ para todo $x\in(a,b)$. Entonces
\[
\int_a^b f(x)\dif x=F(b)-F(a)=\eval{F(x)}{a}{b}.
\]
\end{teorema}

La demostración no exige que $f$ sea continua, solo integrable, y no usa la primera parte. Cuando $f$ es continua también puede deducirse de ella: $G(x)=\int_a^x f$ y $F$ son primitivas de $f$ en $(a,b)$, continuas en $[a,b]$, luego difieren en una constante (\cref{prop:primitivas-constante}), y $F(b)-F(a)=G(b)-G(a)=\int_a^b f$.

\addtocounter{definicion}{1}%

\section{El logaritmo como integral}

El \cref{cor:existencia-primitiva} permite \emph{definir} funciones mediante integrales. El ejemplo más importante es el logaritmo, que puede construirse sin conocer previamente la exponencial~\cite[cap.~18]{spivak}, \cite[cap.~6]{apostol}.

\begin{proposicion}\label{prop:logaritmo}
Sea $L(x)=\int_1^x\frac{\dif t}{t}$ para $x>0$. Entonces $L(1)=0$, $L'(x)=\frac1x$, $L$ es estrictamente creciente y $L(ab)=L(a)+L(b)$ para todos $a,b>0$.
\end{proposicion}

$L$ es, por tanto, el logaritmo natural, $L=\ln$, y su inversa es la exponencial. De la misma forma se definen funciones que no tienen primitiva expresable con funciones elementales, como la función de error $\operatorname{erf}(x)=\frac{2}{\sqrt\pi}\int_0^x e^{-t^{2}}\dif t$ o la función de distribución de la normal estándar, $\Phi(x)=\frac{1}{\sqrt{2\pi}}\int_{-\infty}^{x}e^{-t^{2}/2}\dif t$, que es una integral impropia (\cref{tema:4}). Por el \cref{teo:tfc1}, $\Phi'(x)=\frac{1}{\sqrt{2\pi}}e^{-x^{2}/2}$ aunque $\Phi$ no tenga fórmula elemental.

\section{Áreas entre curvas}

\begin{definicion}[Área entre dos curvas]\label{def:area}
Sean $f,g$ continuas en $[a,b]$. El área de la región comprendida entre sus gráficas y las rectas $x=a$, $x=b$ es $\int_a^b\abs{f(x)-g(x)}\dif x$. Si $f\ge g$ en $[a,b]$, es $\int_a^b(f-g)$.
\end{definicion}

Para calcularla se hallan los puntos donde $f-g$ cambia de signo, se divide $[a,b]$ en subintervalos donde el signo es constante y se suman las integrales de $f-g$ o de $g-f$ según corresponda.

\addtocounter{definicion}{1}%

\addtocounter{definicion}{1}%

\section{Valor medio y simetrías}

\addtocounter{definicion}{1}%

\begin{proposicion}[Integrales de funciones pares e impares]\label{prop:par-impar}
Sea $f$ continua en $[-a,a]$. Si $f$ es par, $\int_{-a}^{a}f=2\int_0^{a}f$; si $f$ es impar, $\int_{-a}^{a}f=0$.
\end{proposicion}

\begin{observacion}[Medias ponderadas]
Si $w\ge0$ es continua con $\int_a^b w=1$, el número $\int_a^b f(x)\,w(x)\dif x$ es una media de $f$ ponderada por $w$. Cuando $w$ es la densidad de una variable aleatoria $X$, es la esperanza $\E[f(X)]$ del módulo~5; el valor medio de la \cref{def:valor-medio} corresponde a la densidad uniforme $w=\frac{1}{b-a}$. El \cref{eje:2-tvm-ponderado} generaliza el \cref{teo:tvm-integral} a este caso.
\end{observacion}

\addtocounter{definicion}{1}%

\addtocounter{definicion}{1}%


% =====================================================================
\chapter{Técnicas de integración}
\label{tema:3}
% =====================================================================

La regla de Barrow reduce el cálculo de $\int_a^b f$ a encontrar una primitiva de $f$. A diferencia de la derivación, que sigue reglas mecánicas, la búsqueda de primitivas es un arte con unas pocas herramientas: el cambio de variable, que invierte la regla de la cadena; la integración por partes, que invierte la regla del producto; la descomposición en fracciones simples, que resuelve por completo las funciones racionales; y las identidades trigonométricas. Algunas funciones elementales, como $e^{-x^{2}}$, no tienen primitiva elemental, y para ellas se recurre a las funciones definidas por integrales del \cref{tema:2} o a las series del \cref{tema:6}. Se siguen el capítulo~19 de Spivak~\cite{spivak}, las §5.5 y~5.6 de~\cite{openstax1} y §3.1 a~3.5 de~\cite{openstax2}, el capítulo~7 de Strang~\cite{strang} y la unidad~4 del curso 18.01SC~\cite{mit1801}.

\section{Cambio de variable}

\begin{teorema}[Cambio de variable]\label{teo:sustitucion}
Sea $\varphi\colon[\alpha,\beta]\to\R$ con derivada continua, y sea $f$ continua en un intervalo $I$ que contiene a $\varphi([\alpha,\beta])$. Entonces
\[
\int_{\alpha}^{\beta}f\bigl(\varphi(t)\bigr)\,\varphi'(t)\dif t=\int_{\varphi(\alpha)}^{\varphi(\beta)}f(x)\dif x .
\]
En forma indefinida: si $F'=f$, entonces $\int f(\varphi(t))\,\varphi'(t)\dif t=F(\varphi(t))+C$.
\end{teorema}

En la práctica se escribe $u=\varphi(t)$, $\dif u=\varphi'(t)\dif t$, y se sustituye formalmente. En una integral definida conviene cambiar también los límites, lo que evita deshacer el cambio.

\addtocounter{definicion}{1}%

El teorema también se usa en sentido inverso: para calcular $\int_a^b f(x)\dif x$ se escribe $x=\varphi(t)$ con $\varphi$ biyectiva de $[\alpha,\beta]$ en $[a,b]$, $\varphi(\alpha)=a$, $\varphi(\beta)=b$, y se integra en $t$.

\addtocounter{definicion}{1}%

\section{Integración por partes}

\begin{teorema}[Integración por partes]\label{teo:partes}
Sean $u,v$ con derivada continua en $[a,b]$. Entonces
\[
\int_a^b u(x)\,v'(x)\dif x=\eval{u(x)\,v(x)}{a}{b}-\int_a^b u'(x)\,v(x)\dif x ,
\]
y en forma indefinida, $\int u\,v'=u\,v-\int u'\,v$, o abreviadamente $\int u\dif v=uv-\int v\dif u$.
\end{teorema}

La técnica es útil cuando $u$ se simplifica al derivarla (un polinomio, un logaritmo, una función trigonométrica inversa) y $v'$ no se complica al integrarla (una exponencial, un seno, un coseno).

\addtocounter{definicion}{1}%

La integración por partes repetida conduce a \emph{fórmulas de reducción}, que expresan una integral con un parámetro $n$ en función de la misma integral con $n-1$ o $n-2$ y se resuelven por inducción (\cref{eje:3-reduccion,eje:3-wallis}).

\section{Funciones racionales y fracciones simples}

Una función racional $P/Q$ es \emph{propia} si $\operatorname{gr}P<\operatorname{gr}Q$. Toda función racional es un polinomio más una función racional propia, por la división de polinomios. Todo polinomio real se factoriza en factores lineales $(x-r)$ y cuadráticos irreducibles $(x^{2}+px+q)$, con $p^{2}-4q<0$.

\begin{teorema}[Descomposición en fracciones simples]\label{teo:fracciones}
Sea $P/Q$ una función racional propia y supóngase $Q$ mónico factorizado como producto de potencias de factores lineales $(x-r)^{k}$ y cuadráticos irreducibles $(x^{2}+px+q)^{m}$. Entonces $P/Q$ es, de forma única, una suma de términos
\[
\frac{A_j}{(x-r)^{j}}\ \ (1\le j\le k)\qquad\text{y}\qquad\frac{B_j\,x+C_j}{(x^{2}+px+q)^{j}}\ \ (1\le j\le m),
\]
uno por cada potencia de cada factor.
\end{teorema}

La demostración general es algebraica~\cite[cap.~19]{spivak}. El caso más frecuente, con raíces reales simples, admite una prueba breve que además da los coeficientes.

\begin{proposicion}[Raíces simples]\label{prop:raices-simples}
Sea $Q(x)=(x-r_1)\cdots(x-r_n)$ con $r_1,\dots,r_n$ distintos y sea $P$ un polinomio de grado menor que $n$. Entonces
\[
\frac{P(x)}{Q(x)}=\sum_{i=1}^{n}\frac{A_i}{x-r_i},\qquad A_i=\frac{P(r_i)}{\prod_{j\neq i}(r_i-r_j)}=\frac{P(r_i)}{Q'(r_i)} .
\]
\end{proposicion}

Las fracciones simples tienen primitivas conocidas: $\int\frac{A\dif x}{x-r}=A\ln\abs{x-r}+C$; $\int\frac{A\dif x}{(x-r)^{j}}=-\frac{A}{(j-1)(x-r)^{j-1}}+C$ si $j\ge2$; y para el factor cuadrático se completa el cuadrado, $x^{2}+px+q=(x+\frac p2)^{2}+s^{2}$ con $s^{2}=q-\frac{p^{2}}{4}>0$, y se separa la parte que es derivada del denominador:
\[
\int\frac{\dif x}{(x+h)^{2}+s^{2}}=\frac1s\arctan\frac{x+h}{s}+C,\qquad\int\frac{2x+p}{x^{2}+px+q}\dif x=\ln(x^{2}+px+q)+C .
\]

\addtocounter{definicion}{1}%

\section{Integrales trigonométricas}

Para $\int\sin^{m}x\cos^{n}x\dif x$ con $m,n\in\N_0$ hay dos casos~\cite[§3.2]{openstax2}.
\begin{enumerate}
  \item Si alguno de los exponentes es impar, por ejemplo $m=2k+1$, se aparta un factor $\sin x$, se escribe $\sin^{2k}x=(1-\cos^{2}x)^{k}$ y se sustituye $u=\cos x$. El caso $n$ impar es simétrico, con $u=\sin x$.
  \item Si ambos son pares, se reduce el grado con $\sin^{2}x=\frac{1-\cos2x}{2}$, $\cos^{2}x=\frac{1+\cos2x}{2}$ y $\sin x\cos x=\frac{\sin2x}{2}$.
\end{enumerate}

\addtocounter{definicion}{1}%

Las raíces cuadráticas se eliminan con \emph{sustituciones trigonométricas}~\cite[§3.3]{openstax2}: $\sqrt{a^{2}-x^{2}}$ con $x=a\sin t$, como en el \cref{ej:circulo}; $\sqrt{a^{2}+x^{2}}$ con $x=a\tan t$; y $\sqrt{x^{2}-a^{2}}$ con $x=a/\cos t$. Por último, la sustitución $t=\tan\frac x2$ transforma cualquier función racional de $\sin x$ y $\cos x$ en una función racional de $t$, pues $\sin x=\frac{2t}{1+t^{2}}$, $\cos x=\frac{1-t^{2}}{1+t^{2}}$ y $\dif x=\frac{2\dif t}{1+t^{2}}$.

\section{Estrategia}

Ante una integral conviene probar, por este orden: (1) simplificar algebraicamente y buscar una primitiva inmediata (\cref{tab:inmediatas}); (2) buscar una función y su derivada en el integrando, para un cambio de variable; (3) si es un producto de un polinomio o logaritmo por una exponencial o función trigonométrica, integrar por partes; (4) si es racional, dividir y descomponer en fracciones simples; (5) si es trigonométrica o contiene raíces de polinomios de grado~2, usar identidades o sustituciones trigonométricas. Toda primitiva se comprueba derivando.

\addtocounter{definicion}{1}%

\addtocounter{definicion}{1}%


% =====================================================================
\chapter{Integrales impropias}
\label{tema:4}
% =====================================================================

La integral de Riemann exige un intervalo cerrado y acotado y una función acotada. Muchas integrales importantes no cumplen alguna de las dos condiciones: la densidad normal está definida en toda la recta real, la función Gamma integra en $[0,\infty)$ y la densidad de una ley de potencia no está acotada cerca del origen si el exponente es negativo. Este tema extiende la integral mediante un paso al límite, estudia cuándo ese límite existe y calcula dos integrales fundamentales para la probabilidad: la función Gamma y la integral gaussiana. Se siguen la §3.7 de~\cite{openstax2}, la §7.5 de Strang~\cite{strang}, el capítulo~10 de Apostol~\cite{apostol} y las sesiones~91 a~93 del curso 18.01SC~\cite{mit1801}.

\section{Integrales en intervalos no acotados}

\begin{definicion}[Integral impropia de primera especie]\label{def:impropia1}
Sea $f\colon[a,\infty)\to\R$ integrable en $[a,b]$ para todo $b>a$. Si existe el límite finito
\[
\int_a^{\infty}f(x)\dif x=\lim_{b\to\infty}\int_a^{b}f(x)\dif x,
\]
se dice que la integral impropia \emph{converge} y que ese es su valor; en otro caso, \emph{diverge}. Se define igual $\int_{-\infty}^{b}f$. Para una función definida en todo $\R$, se dice que $\int_{-\infty}^{\infty}f$ converge si convergen $\int_{-\infty}^{c}f$ y $\int_{c}^{\infty}f$ para algún $c$, y su valor es la suma de ambas.
\end{definicion}

Por la aditividad, la elección de $c$ no afecta ni a la convergencia ni al valor. La definición exige que converjan las dos colas por separado: $\int_{-\infty}^{\infty}x\dif x$ diverge, aunque $\int_{-R}^{R}x\dif x=0$ para todo $R$.

\begin{proposicion}[Integrales $p$ en el infinito]\label{prop:p-infinito}
Sea $p\in\R$. La integral $\int_1^{\infty}\frac{\dif x}{x^{p}}$ converge si y solo si $p>1$, y en ese caso vale $\frac{1}{p-1}$.
\end{proposicion}

\addtocounter{definicion}{1}%

\section{Integrandos no acotados}

\begin{definicion}[Integral impropia de segunda especie]\label{def:impropia2}
Sea $f\colon(a,b]\to\R$ integrable en $[c,b]$ para todo $c\in(a,b)$ y no acotada cerca de $a$. Se define $\int_a^{b}f=\lim_{c\to a^{+}}\int_c^{b}f$ si el límite existe y es finito. Análogamente si la singularidad está en $b$. Si hay una singularidad en un punto interior, o varias, se divide el intervalo y se exige que converja cada trozo.
\end{definicion}

\begin{proposicion}[Integrales $p$ en el origen]\label{prop:p-cero}
La integral $\int_0^{1}\frac{\dif x}{x^{p}}$, con $p>0$, converge si y solo si $p<1$, y en ese caso vale $\frac{1}{1-p}$.
\end{proposicion}

Las \cref{prop:p-infinito,prop:p-cero} son simétricas: $x^{-p}$ es integrable en el infinito si decrece más deprisa que $1/x$, y en el origen si crece más despacio que $1/x$. En particular, $\int_0^{\infty}x^{-p}\dif x$ diverge para todo $p$.

\addtocounter{definicion}{1}%

\section{Criterios de convergencia}

Casi nunca se conoce una primitiva explícita, y conviene decidir la convergencia comparando con integrales conocidas. Los criterios se enuncian en $[a,\infty)$; para las de segunda especie son idénticos.

\begin{teorema}[Criterio de comparación]\label{teo:comparacion-int}
Sean $f,g$ integrables en $[a,b]$ para todo $b>a$, con $0\le f(x)\le g(x)$ para $x\ge a$. Si $\int_a^{\infty}g$ converge, $\int_a^{\infty}f$ converge y $\int_a^{\infty}f\le\int_a^{\infty}g$. Si $\int_a^{\infty}f$ diverge, $\int_a^{\infty}g$ diverge.
\end{teorema}

\begin{teorema}[Comparación en el límite]\label{teo:comparacion-limite}
Sean $f,g>0$ en $[a,\infty)$, integrables en cada $[a,b]$, con $\lim_{x\to\infty}\frac{f(x)}{g(x)}=L\in(0,\infty)$. Entonces $\int_a^{\infty}f$ y $\int_a^{\infty}g$ son ambas convergentes o ambas divergentes.
\end{teorema}

La demostración se pide en el \cref{eje:4-limite-comp}. En la práctica se compara con $x^{-p}$: por ejemplo, $\frac{1}{\sqrt{x^{3}+1}}$ se comporta como $x^{-3/2}$ en el infinito, y $\int_1^{\infty}\frac{\dif x}{\sqrt{x^{3}+1}}$ converge.

\begin{teorema}[Convergencia absoluta]\label{teo:absoluta-int}
Sea $f$ integrable en cada $[a,b]$. Si $\int_a^{\infty}\abs{f}$ converge, también converge $\int_a^{\infty}f$, y $\bigl|\int_a^{\infty}f\bigr|\le\int_a^{\infty}\abs{f}$.
\end{teorema}

\addtocounter{definicion}{1}%

\section{La función Gamma}

Para todo $n\in\N_0$, el límite $\lim_{x\to\infty}x^{n}e^{-x}=0$ (módulo~0.5) expresa que la exponencial domina a cualquier potencia. Gracias a ello las integrales $\int_0^{\infty}x^{n}e^{-x}\dif x$ convergen.

\begin{definicion}[Función Gamma]\label{def:gamma}
Para $s>0$ se define $\Gamma(s)=\int_0^{\infty}x^{s-1}e^{-x}\dif x$.
\end{definicion}

La integral es impropia en el infinito y, si $s<1$, también en el origen. Converge para todo $s>0$: en $(0,1]$, $0<x^{s-1}e^{-x}\le x^{s-1}=x^{-(1-s)}$ con $1-s<1$ (\cref{prop:p-cero}); en $[1,\infty)$, como $x^{s+1}e^{-x}\to0$, existe $x_0$ con $x^{s-1}e^{-x}\le x^{-2}$ para $x\ge x_0$ (\cref{prop:p-infinito}). El \cref{eje:4-gamma-conv} prueba que diverge si $s\le0$.

\begin{proposicion}\label{prop:gamma-factorial}
Para todo $s>0$, $\Gamma(s+1)=s\,\Gamma(s)$. En particular, $\Gamma(1)=1$ y $\Gamma(n+1)=n!$ para todo $n\in\N_0$.
\end{proposicion}

La función Gamma extiende el factorial a valores no enteros; por ejemplo, $\Gamma(\frac12)=\sqrt\pi$ (\cref{eje:4-gamma-medio}).

\section{La integral gaussiana}

La función $e^{-x^{2}}$ no tiene primitiva elemental, pero su integral sobre toda la recta puede calcularse exactamente. Es la constante que normaliza la distribución normal.

\begin{teorema}[Integral gaussiana]\label{teo:gauss}
$\displaystyle\int_{-\infty}^{\infty}e^{-x^{2}}\dif x=\sqrt\pi$.
\end{teorema}

\begin{observacion}[Sobre el paso a coordenadas polares]
La demostración usa dos hechos del cálculo en varias variables: que el producto de dos integrales es una integral doble (teorema de Fubini) y la fórmula del cambio a coordenadas polares, $\dif x\dif y=r\dif r\dif\theta$, de la que el argumento de las coronas es la versión intuitiva. Ambos se justifican en el módulo~0.7; véase también la §14.2 «Changing to Better Coordinates» de Strang~\cite{strang}. El \cref{eje:4-wallis-gauss} obtiene el mismo resultado sin salir de una variable, con las integrales de Wallis del \cref{eje:3-wallis}.
\end{observacion}

\begin{corolario}\label{cor:gauss}
Para $\mu\in\R$ y $\sigma>0$,
\[
\int_{-\infty}^{\infty}e^{-x^{2}/2}\dif x=\sqrt{2\pi},\qquad\int_{-\infty}^{\infty}e^{-(x-\mu)^{2}/(2\sigma^{2})}\dif x=\sigma\sqrt{2\pi}.
\]
En consecuencia, $\varphi(x)=\frac{1}{\sigma\sqrt{2\pi}}e^{-(x-\mu)^{2}/(2\sigma^{2})}$ es una densidad de probabilidad: es positiva e integra~1.
\end{corolario}

\addtocounter{definicion}{1}%

\addtocounter{definicion}{1}%


% =====================================================================
\chapter{Sucesiones y series numéricas}
\label{tema:5}
% =====================================================================

Una serie es una suma de infinitos términos. Como no se pueden sumar infinitos números uno tras otro, la suma se define como el límite de las sumas parciales, y la primera pregunta es si ese límite existe. Este tema repasa los límites de sucesiones, estudia la serie geométrica y la serie armónica, que son los dos modelos de comparación fundamentales, y desarrolla los criterios de convergencia: comparación, cociente, raíz, integral y Leibniz. Se siguen los capítulos~22 y~23 de Spivak~\cite{spivak}, el capítulo~10 de Apostol~\cite{apostol}, las §5.1 a~5.6 de~\cite{openstax2}, las §10.1 a~10.3 de Strang~\cite{strang} y las sesiones~94 y~95 del curso 18.01SC~\cite{mit1801}.

\section{Sucesiones}

\begin{definicion}[Límite de una sucesión]\label{def:limite-sucesion}
Una sucesión $(a_n)_{n\ge1}$ de números reales \emph{converge} a $\ell\in\R$, y se escribe $a_n\to\ell$ o $\lim_{n\to\infty}a_n=\ell$, si para cada $\varepsilon>0$ existe $N\in\N$ tal que $\abs{a_n-\ell}<\varepsilon$ para todo $n\ge N$. Se escribe $a_n\to\infty$ si para cada $M$ existe $N$ con $a_n>M$ para todo $n\ge N$. Una sucesión que no converge \emph{diverge}.
\end{definicion}

El límite, si existe, es único, y los límites de sucesiones cumplen las mismas reglas algebraicas que los de funciones (módulo~0.5): el límite de una suma, de un producto o de un cociente con denominador de límite no nulo es la suma, el producto o el cociente de los límites. También vale el criterio del encaje: si $a_n\le b_n\le c_n$ y $a_n,c_n\to\ell$, entonces $b_n\to\ell$. Si $a_n=f(n)$ y $f(x)\to\ell$ cuando $x\to\infty$, entonces $a_n\to\ell$, lo que permite usar la regla de L'Hôpital.

\begin{teorema}[Convergencia monótona]\label{teo:monotona-acotada}
Toda sucesión creciente y acotada superiormente converge a su supremo. Toda sucesión decreciente y acotada inferiormente converge a su ínfimo.
\end{teorema}

\addtocounter{definicion}{1}%

\section{Series numéricas}

\begin{definicion}[Serie]\label{def:serie}
Dada una sucesión $(a_n)_{n\ge1}$, su \emph{suma parcial} $n$-ésima es $S_n=a_1+\dots+a_n$. La serie $\sum_{n=1}^{\infty}a_n$ \emph{converge} si $(S_n)$ converge, y su \emph{suma} es $\sum_{n=1}^{\infty}a_n=\lim_{n\to\infty}S_n$. En otro caso la serie \emph{diverge}.
\end{definicion}

Si $\sum a_n$ y $\sum b_n$ convergen, también $\sum(\alpha a_n+\beta b_n)$, con suma $\alpha\sum a_n+\beta\sum b_n$, porque las sumas parciales se combinan del mismo modo. Suprimir o añadir un número finito de términos no cambia el carácter de la serie, aunque sí su suma.

\begin{teorema}[Serie geométrica]\label{teo:geometrica}
Sea $r\in\R$. La serie $\sum_{n=0}^{\infty}r^{n}$ converge si y solo si $\abs r<1$, y entonces
\[
\sum_{n=0}^{\infty}r^{n}=\frac{1}{1-r}.
\]
\end{teorema}

\begin{proposicion}[Condición necesaria]\label{prop:condicion-necesaria}
Si $\sum a_n$ converge, entonces $a_n\to0$. Equivalentemente, si $a_n\not\to0$, la serie diverge.
\end{proposicion}

La condición $a_n\to0$ es necesaria pero no suficiente, como muestra el resultado siguiente.

\begin{teorema}[La serie armónica diverge]\label{teo:armonica}
La serie $\sum_{n=1}^{\infty}\frac1n$ diverge; de hecho, $H_{2^{k}}\ge1+\frac k2$ para todo $k\in\N_0$.
\end{teorema}

\addtocounter{definicion}{1}%

\section{Series de términos no negativos}

Si $a_n\ge0$ para todo $n$, las sumas parciales forman una sucesión creciente y, por el \cref{teo:monotona-acotada}, la serie converge si y solo si las sumas parciales están acotadas. De ahí salen todos los criterios de esta sección.

\begin{teorema}[Criterio de comparación]\label{teo:comparacion-series}
Sean $0\le a_n\le b_n$ para todo $n$. Si $\sum b_n$ converge, $\sum a_n$ converge y $\sum a_n\le\sum b_n$. Si $\sum a_n$ diverge, $\sum b_n$ diverge.
\end{teorema}

\begin{teorema}[Comparación en el límite]\label{teo:comparacion-limite-series}
Sean $a_n,b_n>0$ con $\frac{a_n}{b_n}\to L\in(0,\infty)$. Entonces $\sum a_n$ y $\sum b_n$ tienen el mismo carácter.
\end{teorema}

\begin{teorema}[Criterio de la integral]\label{teo:criterio-integral}
Sea $f\colon[1,\infty)\to\R$ positiva y decreciente y sea $a_k=f(k)$. Para todo $n\ge1$,
\[
\int_1^{n+1}f(x)\dif x\le\sum_{k=1}^{n}a_k\le a_1+\int_1^{n}f(x)\dif x .
\]
En consecuencia, $\sum a_k$ converge si y solo si $\int_1^{\infty}f$ converge.
\end{teorema}

\addtocounter{figure}{1}%

\begin{corolario}[Series $p$]\label{cor:series-p}
La serie $\sum_{n=1}^{\infty}\frac{1}{n^{p}}$ converge si y solo si $p>1$. Además, $\ln(n+1)\le H_n\le1+\ln n$ para todo $n\in\N$.
\end{corolario}

Así, la serie armónica diverge, pero muy despacio: $H_{1000}$ está entre $\ln1001\approx6{,}91$ y $1+\ln1000\approx7{,}91$. En cambio, $\sum\frac{1}{n^{2}}$ converge, y su suma es menor que $1+\int_1^\infty x^{-2}\dif x=2$ (su valor exacto, $\frac{\pi^{2}}{6}\approx1{,}645$, se obtiene con otras técnicas).

\section{Criterios del cociente y de la raíz}

\begin{teorema}[Criterio del cociente]\label{teo:cociente}
Sea $a_n>0$ con $\frac{a_{n+1}}{a_n}\to L$. Si $L<1$, $\sum a_n$ converge; si $L>1$ (incluido $L=\infty$), diverge. Si $L=1$, el criterio no decide.
\end{teorema}

\begin{teorema}[Criterio de la raíz]\label{teo:raiz}
Sea $a_n\ge0$ con $\sqrt[n]{a_n}\to L$. Si $L<1$, $\sum a_n$ converge; si $L>1$, diverge. Si $L=1$, el criterio no decide.
\end{teorema}

La demostración del criterio del cociente es análoga y se pide en el \cref{eje:5-cociente-dem}. Ambos criterios comparan con una serie geométrica, y por eso no deciden en el caso $L=1$: tanto $\sum\frac1n$, divergente, como $\sum\frac{1}{n^{2}}$, convergente, tienen $L=1$. Son especialmente útiles con factoriales y potencias $n$-ésimas.

\addtocounter{definicion}{1}%

\section{Series alternadas y convergencia absoluta}

\begin{definicion}[Convergencia absoluta y condicional]\label{def:absoluta}
La serie $\sum a_n$ converge \emph{absolutamente} si $\sum\abs{a_n}$ converge. Si $\sum a_n$ converge pero $\sum\abs{a_n}$ diverge, converge \emph{condicionalmente}.
\end{definicion}

\begin{teorema}\label{teo:absoluta-series}
Toda serie absolutamente convergente es convergente, y $\bigl|\sum a_n\bigr|\le\sum\abs{a_n}$.
\end{teorema}

\begin{teorema}[Criterio de Leibniz]\label{teo:leibniz}
Sea $(b_n)$ decreciente con $b_n\to0$ (luego $b_n\ge0$). La serie alternada $\sum_{n=1}^{\infty}(-1)^{n+1}b_n=b_1-b_2+b_3-\cdots$ converge, y su suma $S$ cumple $\abs{S-S_n}\le b_{n+1}$ para todo $n$.
\end{teorema}

\addtocounter{definicion}{1}%

\begin{observacion}[Reordenaciones]
Una serie absolutamente convergente puede sumarse en cualquier orden con el mismo resultado. Con una serie condicionalmente convergente no ocurre así: un teorema de Riemann afirma que sus términos pueden reordenarse para obtener cualquier suma prefijada~\cite[cap.~23]{spivak}. El ejercicio~25 de la §10.3 de Strang~\cite{strang} muestra una reordenación de la serie armónica alternada cuya suma es $\frac12\ln2$.
\end{observacion}

\addtocounter{definicion}{1}%

\addtocounter{definicion}{1}%


% =====================================================================
\chapter{Series de potencias y series de Taylor}
\label{tema:6}
% =====================================================================

Una serie de potencias es un «polinomio de grado infinito». Dentro de su intervalo de convergencia se comporta como un polinomio: puede derivarse e integrarse término a término. Las funciones elementales $e^{x}$, $\sin x$, $\cos x$, $\ln(1+x)$ y $(1+x)^{a}$ son sumas de series de potencias, sus series de Taylor, y esa representación permite calcularlas con la precisión deseada, obtener límites y definir el número $e$. Este tema estudia el radio de convergencia, las operaciones con series y los desarrollos fundamentales, con demostración de que convergen a la función. Se siguen los capítulos~20 y~24 de Spivak~\cite{spivak}, los capítulos~7 y~11 de Apostol~\cite{apostol}, las §6.1 a~6.4 de~\cite{openstax2}, las §10.4 y~10.5 de Strang~\cite{strang} y las sesiones~97 a~100 del curso 18.01SC~\cite{mit1801}.

\section{Series de potencias y radio de convergencia}

\begin{definicion}[Serie de potencias]\label{def:serie-potencias}
Una \emph{serie de potencias} centrada en $c\in\R$ es una serie de la forma $\sum_{n=0}^{\infty}a_n(x-c)^{n}$, donde $(a_n)$ es una sucesión de números reales y $x$ es una variable real. Su \emph{dominio de convergencia} es el conjunto de los $x$ para los que converge.
\end{definicion}

Toda serie de potencias converge en $x=c$, donde vale $a_0$. El resultado siguiente muestra que el dominio de convergencia es siempre un intervalo centrado en $c$.

\begin{lema}[Lema de Abel]\label{lem:abel}
Si $\sum a_n(x_0-c)^{n}$ converge para algún $x_0\neq c$, entonces $\sum a_n(x-c)^{n}$ converge absolutamente para todo $x$ con $\abs{x-c}<\abs{x_0-c}$.
\end{lema}

\begin{teorema}[Radio de convergencia]\label{teo:radio}
Para cada serie de potencias $\sum a_n(x-c)^{n}$ existe $R\in[0,\infty]$, su \emph{radio de convergencia}, tal que la serie converge absolutamente si $\abs{x-c}<R$ y diverge si $\abs{x-c}>R$. Si existe $\lim_{n\to\infty}\Bigl|\frac{a_n}{a_{n+1}}\Bigr|$ (finito o infinito), ese límite es $R$.
\end{teorema}

El intervalo $(c-R,c+R)$ es el \emph{intervalo de convergencia}. En los extremos $x=c\pm R$ el teorema no dice nada, y cada caso se estudia aparte con los criterios del \cref{tema:5}.

\addtocounter{definicion}{1}%

\section{Derivación e integración término a término}

\begin{teorema}[Derivación e integración de series de potencias]\label{teo:termino-a-termino}
Sea $f(x)=\sum_{n=0}^{\infty}a_n(x-c)^{n}$ con radio de convergencia $R>0$. Las series
\[
\sum_{n=1}^{\infty}n\,a_n(x-c)^{n-1}\qquad\text{y}\qquad\sum_{n=0}^{\infty}\frac{a_n}{n+1}(x-c)^{n+1}
\]
tienen el mismo radio $R$; $f$ es derivable en $(c-R,c+R)$, con $f'$ igual a la primera serie, y $\int_c^{x}f(t)\dif t$ es igual a la segunda para todo $x$ de ese intervalo.
\end{teorema}

La demostración usa la convergencia uniforme de la serie en los intervalos cerrados contenidos en $(c-R,c+R)$ y puede consultarse en~\cite[cap.~24]{spivak} o \cite[cap.~11]{apostol}. Aplicando el teorema repetidamente, $f$ tiene derivadas de todos los órdenes en el intervalo de convergencia.

\begin{corolario}[Unicidad de los coeficientes]\label{cor:coeficientes}
Si $f(x)=\sum_{n=0}^{\infty}a_n(x-c)^{n}$ con $R>0$, entonces $a_n=\frac{f^{(n)}(c)}{n!}$ para todo $n$. En particular, dos series de potencias que coinciden en un entorno de $c$ tienen los mismos coeficientes.
\end{corolario}

\addtocounter{definicion}{1}%

\section{Series de Taylor}

\begin{definicion}[Serie de Taylor]\label{def:serie-taylor}
Sea $f$ con derivadas de todos los órdenes en un intervalo abierto que contiene a $c$. Su \emph{serie de Taylor} en $c$ es $\sum_{n=0}^{\infty}\frac{f^{(n)}(c)}{n!}(x-c)^{n}$. Para $c=0$ se llama también \emph{serie de Maclaurin}. La suma parcial de grado $n$ es el polinomio de Taylor $T_nf(x)$, y $R_nf(x)=f(x)-T_nf(x)$ es el \emph{resto}.
\end{definicion}

Por el \cref{cor:coeficientes}, si $f$ es suma de alguna serie de potencias centrada en $c$, esa serie es su serie de Taylor. Pero la serie de Taylor de $f$ puede converger y no tener por suma $f$ (\cref{eje:6-no-analitica}). La herramienta para decidirlo es el teorema de Taylor del módulo~0.5.

\begin{teorema}[Taylor con resto de Lagrange]\label{teo:taylor-lagrange}
Sea $f$ con $n+1$ derivadas en un intervalo abierto $I$ que contiene a $c$. Para cada $x\in I$ existe $\xi$ entre $c$ y $x$ tal que
\[
R_nf(x)=\frac{f^{(n+1)}(\xi)}{(n+1)!}(x-c)^{n+1}.
\]
\end{teorema}

\begin{proposicion}\label{prop:taylor-converge}
La serie de Taylor de $f$ en $c$ converge a $f(x)$ si y solo si $R_nf(x)\to0$ cuando $n\to\infty$.
\end{proposicion}

\section{Los desarrollos fundamentales}

\begin{teorema}[Exponencial]\label{teo:exp}
Para todo $x\in\R$, $\displaystyle e^{x}=\sum_{n=0}^{\infty}\frac{x^{n}}{n!}=1+x+\frac{x^{2}}{2}+\frac{x^{3}}{6}+\cdots$.
\end{teorema}

Del mismo modo, como las derivadas de $\sin$ y $\cos$ son $\pm\sin$ y $\pm\cos$, acotadas por~1, sus restos cumplen $\abs{R_n(x)}\le\frac{\abs x^{n+1}}{(n+1)!}\to0$, y para todo $x\in\R$ (\cref{eje:6-seno-dem})
\[
\sin x=\sum_{n=0}^{\infty}\frac{(-1)^{n}x^{2n+1}}{(2n+1)!}=x-\frac{x^{3}}{6}+\cdots,\qquad
\cos x=\sum_{n=0}^{\infty}\frac{(-1)^{n}x^{2n}}{(2n)!}=1-\frac{x^{2}}{2}+\cdots.
\]

\begin{teorema}[Logaritmo]\label{teo:ln}
Para todo $x\in(-1,1]$, $\displaystyle\ln(1+x)=\sum_{n=1}^{\infty}\frac{(-1)^{n+1}}{n}x^{n}$. En particular, $\displaystyle\ln2=1-\frac12+\frac13-\frac14+\cdots$.
\end{teorema}

La demostración cubre el extremo $x=1$, que el \cref{teo:termino-a-termino} no alcanza, y da además una cota explícita del error. La serie no sirve para $x>1$, porque allí su radio de convergencia, $R=1$, ya no la hace converger, aunque $\ln(1+x)$ esté definido.

\begin{teorema}[Serie binomial]\label{teo:binomial}
Sea $a\in\R$. Para todo $x\in(-1,1)$,
\[
(1+x)^{a}=\sum_{n=0}^{\infty}\binom{a}{n}x^{n},\qquad\binom{a}{n}=\frac{a(a-1)\cdots(a-n+1)}{n!},\quad\binom{a}{0}=1 .
\]
Si $a\in\N_0$, la serie es finita y se reduce al binomio de Newton.
\end{teorema}

Si $a\notin\N_0$, ningún coeficiente es nulo, $\bigl|\binom{a}{n}\big/\binom{a}{n+1}\bigr|=\frac{n+1}{\abs{a-n}}\to1$ y el radio es~1. La demostración de que la suma es $(1+x)^{a}$ se pide en el \cref{eje:6-binomial-dem}. Por ejemplo, con $a=\frac12$: $\sqrt{1+x}=1+\frac x2-\frac{x^{2}}{8}+\frac{x^{3}}{16}-\cdots$.

\begin{table}[h]
\centering
\small
\renewcommand{\arraystretch}{1.25}
\begin{tabular}{@{}>{$\displaystyle}l<{$}>{$\displaystyle}l<{$}>{$}l<{$}@{}}
\toprule
\text{Función} & \text{Serie de Maclaurin} & \text{Validez} \\
\midrule
\frac{1}{1-x} & \sum_{n\ge0}x^{n} & (-1,1) \\
e^{x} & \sum_{n\ge0}\frac{x^{n}}{n!} & \R \\
\sin x & \sum_{n\ge0}\frac{(-1)^{n}x^{2n+1}}{(2n+1)!} & \R \\
\cos x & \sum_{n\ge0}\frac{(-1)^{n}x^{2n}}{(2n)!} & \R \\
\ln(1+x) & \sum_{n\ge1}\frac{(-1)^{n+1}x^{n}}{n} & (-1,1] \\
\arctan x & \sum_{n\ge0}\frac{(-1)^{n}x^{2n+1}}{2n+1} & [-1,1] \\
(1+x)^{a} & \sum_{n\ge0}\binom{a}{n}x^{n} & (-1,1) \\
\bottomrule
\end{tabular}
\caption{Desarrollos fundamentales. La validez en $x=\pm1$ de la serie del arcotangente se prueba en el \cref{eje:6-arctan}.}
\label{tab:desarrollos}
\end{table}

\section{El número \texorpdfstring{$e$}{e}}

El número $e$ es la base de la exponencial, $e=e^{1}$. Por el \cref{teo:exp},
\[
e=\sum_{n=0}^{\infty}\frac{1}{n!}=1+1+\frac12+\frac16+\frac{1}{24}+\cdots\approx2{,}718281828 .
\]

\begin{proposicion}\label{prop:e-limite}
$\displaystyle\lim_{n\to\infty}\Bigl(1+\frac1n\Bigr)^{n}=e$.
\end{proposicion}

\begin{teorema}[El número $e$ es irracional]\label{teo:e-irracional}
Sea $s_q=\sum_{k=0}^{q}\frac{1}{k!}$. Para todo $q\in\N$, $0<e-s_q<\frac{1}{q!\,q}$. En consecuencia, $e\notin\Q$.
\end{teorema}

La cota del teorema también es práctica: con $q=10$, $s_{10}$ aproxima $e$ con error menor que $\frac{1}{10!\cdot10}\approx2{,}8\cdot10^{-8}$.

\addtocounter{definicion}{1}%

\addtocounter{definicion}{1}%



\backmatter
% =====================================================================
\begin{thebibliography}{99}
\addcontentsline{toc}{chapter}{Bibliografía}
\markboth{Bibliografía}{}
\small

\bibitem{apostol}
T.~M. Apostol.
\newblock \emph{Calculus}, vol.~1: \emph{One-Variable Calculus, with an Introduction to Linear Algebra}, 2.ª ed.
\newblock John Wiley \& Sons, 1967.
\newblock Caps.~1 «The Concepts of Integral Calculus», 5 «The Relation between Integration and Differentiation», 7 «Polynomial Approximations to Functions», 10 «Sequences, Infinite Series, and Improper Integrals» y 11 «Sequences and Series of Functions».

\bibitem{mit1801}
D.~Jerison \emph{et al.}
\newblock 18.01SC: Single Variable Calculus.
\newblock Massachusetts Institute of Technology, MIT OpenCourseWare, otoño de 2010. Unidad~3 «The Definite Integral and its Applications», unidad~4 «Techniques of Integration» y unidad~5 «Exploring the Infinite»; colecciones de ejercicios \emph{18.01 Exercises}, unidad~3 «Integration» (secciones 3A--3G) y unidad~5 «Integration Techniques» (secciones 5A--5F).
\newblock \url{https://ocw.mit.edu/courses/18-01sc-single-variable-calculus-fall-2010/}

\bibitem{openstax1}
G.~Strang y E.~Herman.
\newblock \emph{Calculus Volume~1}.
\newblock OpenStax, Rice University, 2016. Cap.~5 «Integration» (§5.1--5.7).
\newblock \url{https://openstax.org/details/books/calculus-volume-1}

\bibitem{openstax2}
G.~Strang y E.~Herman.
\newblock \emph{Calculus Volume~2}.
\newblock OpenStax, Rice University, 2016. Caps.~3 «Techniques of Integration» (§3.1--3.7), 5 «Sequences and Series» (§5.1--5.6) y 6 «Power Series» (§6.1--6.4).
\newblock \url{https://openstax.org/details/books/calculus-volume-2}

\bibitem{spivak}
M.~Spivak.
\newblock \emph{Calculus}, 4.ª ed.
\newblock Publish or Perish, 2008.
\newblock Caps.~13 «Integrals», 14 «The Fundamental Theorem of Calculus», 18 «The Logarithm and Exponential Functions», 19 «Integration in Elementary Terms», 20 «Approximation by Polynomial Functions», 22 «Infinite Sequences», 23 «Infinite Series» y 24 «Uniform Convergence and Power Series».

\bibitem{stewart}
J.~Stewart.
\newblock \emph{Calculus: Early Transcendentals}, 8.ª ed.
\newblock Cengage Learning, 2015.

\bibitem{strang}
G.~Strang.
\newblock \emph{Calculus}, 3.ª ed.
\newblock Wellesley-Cambridge Press, 2017. Edición en línea en MIT OpenCourseWare, RES.18-001, otoño de 2023. Caps.~5 «Integrals», 7 «Techniques of Integration», 10 «Infinite Series» y §14.2 «Changing to Better Coordinates».
\newblock \url{https://ocw.mit.edu/courses/res-18-001-calculus-fall-2023/}

\end{thebibliography}


\end{document}
